% ==============================================================================
% ISTITUZIONI DI ALGEBRA
% Teoria Rigorosa, Dimostrazioni, Eserciziario e Guida all'Esame
%
% Basato sui corsi, monografie e registri ufficiali del Prof. Giovanni Cutolo
% Dipartimento di Matematica e Applicazioni "Renato Caccioppoli"
% Università degli Studi di Napoli Federico II
%
% Curatore: Fagghino
% Anno Accademico: 2025/2026
% ==============================================================================

\documentclass[11pt,a4paper,twoside,openright]{book}

% ==============================================================================
% PREAMBLE - ISTITUZIONI DI ALGEBRA (PROF. GIOVANNI CUTOLO)
% Corso di Laurea in Informatica - Università degli Studi di Napoli Federico II
% Curatore: Fagghino
% ==============================================================================

\usepackage[utf8]{inputenc}
\usepackage[T1]{fontenc}
\usepackage[italian]{babel}

% Geometria pagina e margini
\usepackage[a4paper,top=2.5cm,bottom=2.5cm,left=2.5cm,right=2.5cm,headheight=14.5pt]{geometry}

% Matematica e simboli
\usepackage{amsmath,amssymb,amsthm,mathtools,amsfonts,mathrsfs}
\usepackage{bm}

% Grafica e diagrammi
\usepackage{tikz}
\usepackage{tikz-cd}
\usetikzlibrary{shapes,arrows,positioning,calc,matrix,fit,backgrounds}

% Box decorati per definizioni, teoremi ed osservazioni
\usepackage[most]{tcolorbox}

% Gestione elenchi
\usepackage{enumitem}
\setlist{nosep}

% Intestazioni e piè di pagina
\usepackage{fancyhdr}
\pagestyle{fancy}
\fancyhf{}
\fancyhead[LE,RO]{\thepage}
\fancyhead[RE]{\nouppercase{\leftmark}}
\fancyhead[LO]{\nouppercase{\rightmark}}
\renewcommand{\headrulewidth}{0.4pt}

% Collegamenti ipertestuali e riferimenti
\usepackage{xcolor}
\definecolor{UniNaBlue}{RGB}{0, 48, 135}
\definecolor{CutoloRed}{RGB}{160, 20, 20}
\definecolor{DefGreen}{RGB}{20, 110, 60}
\definecolor{ExOrange}{RGB}{180, 80, 0}
\definecolor{NoteGray}{RGB}{80, 80, 90}

\usepackage{hyperref}
\hypersetup{
    colorlinks=true,
    linkcolor=UniNaBlue,
    citecolor=CutoloRed,
    urlcolor=UniNaBlue,
    pdftitle={Istituzioni di Algebra - Appunti Ufficiali Prof. Cutolo},
    pdfauthor={Fagghino},
    pdfsubject={Corso di Algebra - Università degli Studi di Napoli Federico II}
}

\usepackage{cleveref}

% ==============================================================================
% MACRO MATEMATICHE E NOTAZIONI UFFICIALI CUTOLO
% ==============================================================================

% Insiemi numerici standard
\newcommand{\N}{\mathbb{N}}
\newcommand{\Z}{\mathbb{Z}}
\newcommand{\Q}{\mathbb{Q}}
\newcommand{\R}{\mathbb{R}}
\newcommand{\C}{\mathbb{C}}
\newcommand{\K}{\mathbb{K}}
\newcommand{\F}{\mathbb{F}}

% Insieme delle parti
\renewcommand{\P}{\mathscr{P}}
\newcommand{\Pfin}{\mathscr{P}_{\text{fin}}}

% Relazioni e connettivi speciali
\newcommand{\vuoto}{\varnothing}
\newcommand{\smin}{\setminus}
\newcommand{\implica}{\Rightarrow}
\newcommand{\Implica}{\Longrightarrow}
\newcommand{\shiff}{\Leftrightarrow}
\newcommand{\Shiff}{\Longleftrightarrow}
\newcommand{\sseq}{\subseteq}
\newcommand{\ds}{\mathbin{\scriptstyle\triangle}} % Differenza simmetrica
\newcommand{\xor}{\mathbin{\mathsf{XOR}}}
\newcommand{\nor}{\mathbin{\mathsf{NOR}}}
\newcommand{\nand}{\mathbin{\mathsf{NAND}}}

% Notazione specifica del Prof. Cutolo per immagini e controimmagini
% Freccia verso destra per l'immagine f(X)
\newcommand{\vecf}{\vec{f}}
% Freccia verso sinistra per la controimmagine / anti-immagine (NON f^-1!)
\newcommand{\antivec}[2]{#1^{\raise #2pt\hbox{$\!\!\scriptstyle\leftarrow\!\!$}}}
\newcommand{\antivecf}{\antivec{f}{3}}
\newcommand{\antivecg}{\antivec{g}{2}}

% Operatori e insiemi di strutture
\DeclareMathOperator{\id}{id}
\DeclareMathOperator{\im}{im}
\DeclareMathOperator{\Ker}{Ker}
\DeclareMathOperator{\Sym}{Sym}
\DeclareMathOperator{\Aut}{Aut}
\DeclareMathOperator{\Hom}{Hom}
\DeclareMathOperator{\End}{End}
\DeclareMathOperator{\Eq}{Eq}
\DeclareMathOperator{\Part}{Part}
\DeclareMathOperator{\Rel}{Rel}
\DeclareMathOperator{\Corr}{Corr}
\DeclareMathOperator{\Map}{Map}
\DeclareMathOperator{\car}{char}
\DeclareMathOperator{\mcd}{MCD}
\DeclareMathOperator{\mcm}{mcm}
\DeclareMathOperator{\minor}{Minor}
\DeclareMathOperator{\maggior}{Maggior}
\DeclareMathOperator{\degpoly}{deg}

% Generatori e classi resto
\newcommand{\gen}[1]{\langle #1 \rangle}
\newcommand{\classe}[1]{[#1]}
\newcommand{\modulo}[1]{\ (\mathrm{mod}\ #1)}

% ==============================================================================
% AMBIENTI TCOLORBOX (Definizioni, Teoremi, Proposizioni, Esercizi)
% ==============================================================================

% Stile per Definizioni
\newtcbtheorem[number within=chapter]{definizione}{Definizione}{
    enhanced,
    colback=DefGreen!5!white,
    colframe=DefGreen!80!black,
    fonttitle=\bfseries,
    separator sign dash,
    boxrule=0.8pt,
    arc=3pt,
    top=6pt,bottom=6pt,left=8pt,right=8pt
}{def}

% Stile per Teoremi
\newtcbtheorem[use counter from=definizione]{teorema}{Teorema}{
    enhanced,
    colback=UniNaBlue!5!white,
    colframe=UniNaBlue!90!black,
    fonttitle=\bfseries,
    separator sign dash,
    boxrule=0.8pt,
    arc=3pt,
    top=6pt,bottom=6pt,left=8pt,right=8pt
}{teo}

% Stile per Proposizioni / Lemmi
\newtcbtheorem[use counter from=definizione]{proposizione}{Proposizione}{
    enhanced,
    colback=UniNaBlue!3!white,
    colframe=UniNaBlue!60!black,
    fonttitle=\bfseries,
    separator sign dash,
    boxrule=0.6pt,
    arc=2pt,
    top=5pt,bottom=5pt,left=7pt,right=7pt
}{prop}

% Stile per Esercizi e Problemi Svolti
\newtcbtheorem[number within=chapter]{esercizio}{Esercizio d'Esame}{
    enhanced,
    colback=ExOrange!4!white,
    colframe=ExOrange!85!black,
    fonttitle=\bfseries,
    separator sign dash,
    boxrule=0.8pt,
    arc=3pt,
    top=6pt,bottom=6pt,left=8pt,right=8pt
}{es}

% Box per i Tranelli dell'Orale e le Osservazioni di Cutolo
\newtcolorbox{avvisoCutolo}[1][]{
    enhanced,
    colback=CutoloRed!5!white,
    colframe=CutoloRed!85!black,
    fonttitle=\bfseries,
    title={\textbf{All'Orale con Cutolo: Attenzione al Rigore!}},
    boxrule=0.8pt,
    arc=3pt,
    top=6pt,bottom=6pt,left=8pt,right=8pt,
    #1
}

% Box per l'intuizione o la panoramica
\newtcolorbox{intuizione}[1][]{
    enhanced,
    colback=gray!5!white,
    colframe=gray!60!black,
    fonttitle=\bfseries,
    title={Intuizione e Quadro Generale},
    boxrule=0.6pt,
    arc=2pt,
    top=5pt,bottom=5pt,left=7pt,right=7pt,
    #1
}


\begin{document}

% ------------------------------------------------------------------------------
% FRONTESPIZIO
% ------------------------------------------------------------------------------
\begin{titlepage}
    \centering
    \vspace*{1.5cm}
    
    {\scshape\LARGE Università degli Studi di Napoli Federico II \par}
    \vspace{0.3cm}
    {\scshape\Large Corso di Laurea Triennale in Informatica \par}
    \vspace{0.2cm}
    {\large Dipartimento di Matematica e Applicazioni ``Renato Caccioppoli'' \par}
    
    \vspace{1.8cm}
    \rule{\linewidth}{1.2pt} \\[0.4cm]
    {\Huge \bfseries ISTITUZIONI DI ALGEBRA \par}
    \vspace{0.3cm}
    {\Large \bfseries Teoria, Dimostrazioni Complete, Esercizi e Guida all'Esame \par}
    \rule{\linewidth}{1.2pt} \\[0.5cm]
    
    {\large \itshape Dalle lezioni frontali, monografie DMA e registri ufficiali del \par}
    \vspace{0.2cm}
    {\Large \bfseries Prof. Giovanni Cutolo \par}
    
    \vfill
    
    \begin{minipage}{0.48\textwidth}
        \begin{flushleft} \large
            \emph{Insegnamento:}\\
            \textbf{Algebra (LT Informatica)}\\
            Codice: \textbf{U6282} (6 CFU)\\
            Canale 2 (a.a. 2025/2026)
        \end{flushleft}
    \end{minipage}
    \hfill
    \begin{minipage}{0.48\textwidth}
        \begin{flushright} \large
            \emph{A cura di:}\\
            \textbf{Fagghino}\\
            \href{https://github.com/UniNa-Fagghino/Algebra-Cutolo}{\texttt{UniNa-Fagghino}}
        \end{flushright}
    \end{minipage}
    
    \vspace{1.5cm}
    {\large Versione Ufficiale Integrata e Verificata \par}
    {\large Anno Accademico 2025--2026 \par}
\end{titlepage}

% ------------------------------------------------------------------------------
% PREFAZIONE E INDICE
% ------------------------------------------------------------------------------
\frontmatter

\chapter*{Prefazione e Metodologia di Studio}
\addcontentsline{toc}{chapter}{Prefazione e Metodologia di Studio}

\section*{Origine e Scopo di Questa Opera}
Questo testo nasce con un obiettivo preciso e ambizioso: fornire allo studente del Corso di Laurea Triennale in Informatica dell'Università degli Studi di Napoli Federico II **il trattato più completo, rigoroso e autosufficiente possibile** per la preparazione dell'esame di \textbf{Algebra} tenuto dal \textbf{Prof. Giovanni Cutolo}.

La disciplina dell'Algebra Astratta rappresenta spesso il primo vero impatto dello studente universitario con la matematica pura e formale. A differenza dell'analisi o del calcolo numerico, l'algebra non si focalizza sul ``calcolare numeri'', bensì sull'identificare \emph{strutture}, \emph{invarianti}, \emph{simmetrie} e \emph{relazioni formali}.

Il corso del Prof. Cutolo è noto per il suo elevatissimo grado di rigore: ogni affermazione dev'essere motivata a partire dagli assiomi, ogni quantificatore deve trovarsi al posto giusto e ogni dimostrazione deve possedere una struttura logica ineccepibile. Questa dispensa raccoglie, unifica e armonizza:
\begin{enumerate}
    \item Tutte le \textbf{monografie ufficiali} pubblicate dal docente sul portale del Dipartimento di Matematica ed Applicazioni (DMA) ``Renato Caccioppoli''.
    \item Il \textbf{registro analitico delle lezioni} svolte giorno per giorno per l'a.a. 2025/2026.
    \item Il patrimonio di note, appunti manoscritti e risoluzioni di prove scritte condivisi negli anni dalla comunità studentesca della Federico II.
    \item Il censimento dettagliato di tutte le \textbf{domande d'orale ricorrenti}, dei trabocchetti terminologici e degli errori concettuali più frequenti rilevati nelle sessioni d'esame.
\end{enumerate}

\section*{Le Convenzioni Notazionali di Cutolo}
Nel corso del testo si adotta rigorosamente la notazione ufficiale utilizzata dal docente a lezione e nelle prove scritte. Si evidenziano fin da ora alcune convenzioni cardine:
\begin{itemize}
    \item \textbf{I Numeri Naturali $\N$:} per convenzione fissata dal docente, l'insieme dei numeri naturali \textbf{include lo zero}:
    \[
        \N = \{0, 1, 2, 3, \dots\}
    \]
    I naturali strettamente positivi vengono denotati con $\N^* = \N \smin \{0\} = \{1, 2, 3, \dots\}$.
    \item \textbf{Immagine e Controimmagine ($\vecf$ vs $\antivecf$):} Sia $f: A \to B$ un'applicazione. Per indicare l'immagine di un sottoinsieme $X \sseq A$, si usa la notazione con freccia in avanti:
    \[
        \vecf(X) = \{f(x) \in B \mid x \in X\}
    \]
    Per la controimmagine (o anti-immagine) di un sottoinsieme $Y \sseq B$, il Prof. Cutolo utilizza sistematicamente la freccia all'indietro:
    \[
        \antivecf(Y) = \{x \in A \mid f(x) \in Y\}
    \]
    \begin{avvisoCutolo}
        \textbf{Mai usare $f^{-1}(Y)$ all'orale se $f$ non è biettiva!} Il docente ribadisce costantemente che scrivere $f^{-1}$ evoca l'applicazione inversa, che \emph{non esiste} se $f$ non è biettiva. L'operazione di controimmagine $\antivecf: \P(B) \to \P(A)$ è invece sempre definita per qualunque applicazione.
    \end{avvisoCutolo}
    \item \textbf{Composizione di Applicazioni (Prodotto Operativo):} Dati $f: A \to B$ e $g: B \to C$, l'applicazione composta si indica con $g \circ f$ (o semplicemente $gf$), ove $(g \circ f)(a) = g(f(a))$.
    \item \textbf{Relazioni d'Equivalenza e Quozienti:} La classe d'equivalenza di un elemento $a$ modulo una relazione $\mathcal{R}$ si denota con $[a]_\mathcal{R}$ o $[a]$. L'insieme quoziente è $A/\mathcal{R}$. La proiezione canonica è $\pi_\mathcal{R}: a \mapsto [a]_\mathcal{R}$.
\end{itemize}

\section*{Struttura dell'Esame e Come Usare Questa Dispensa}
L'esame di Algebra (codice U6282, 6 CFU) per il Corso di Laurea in Informatica prevede:
\begin{enumerate}
    \item \textbf{Prova Scritta:} della durata di 2 ore (secondo il regolamento aggiornato del nuovo ordinamento), articolata in una prima sezione di quesiti a risposta sintetica/teorica su definizioni e proprietà, seguita da esercizi applicativi (studio di una relazione o poset con diagramma di Hasse, equazioni congruenziali lineari con algoritmo di Bézout o sistemi cinesi, decomposizione di polinomi o studio di una struttura algebrica).
    \item \textbf{Prova Orale:} colloquio individuale approfondito. Il punto di partenza può essere la discussione del compito scritto, ma la verifica verte sistematicamente sulla padronanza formale delle definizioni, degli enunciati dei teoremi e delle rispettive dimostrazioni.
\end{enumerate}

Si raccomanda vivamente di non limitarsi alla risoluzione meccanica degli esercizi: ogni passaggio di calcolo (ad esempio nella risoluzione di un'equazione diofantea o nella determinazione degli elementi invertibili di un anello) poggia su un solido teorema algebrico che all'orale vi verrà puntualmente chiesto di enunciare e dimostrare.


\tableofcontents

% ------------------------------------------------------------------------------
% PARTE I: FONDAMENTI, INSIEMI E COMBINATORIA
% ------------------------------------------------------------------------------
\mainmatter

\part{Logica, Insiemi e Fondamenti del Calcolo}
\chapter{Logica Proposizionale e Calcolo dei Predicati}
\label{chap:logica}

La logica matematica costituisce il linguaggio fondamentale su cui poggia l'intero edificio dell'algebra moderna. In questo capitolo introduciamo gli strumenti del calcolo proposizionale e del calcolo dei predicati, indispensabili per formulare definizioni rigorose ed elaborare dimostrazioni formali ineccepibili.

\section{Proposizioni e Valori di Verità}

\begin{definizione}{Proposizione}{proposizione}
Una \textbf{proposizione} (o enunciato) è una frase dichiarativa per la quale ha senso affermare in modo univoco se sia \emph{vera} oppure \emph{falsa}.

Convenzionalmente, denotiamo i due \textbf{valori di verità} con i simboli:
\[
    1 \quad (\text{Vero}, \top), \qquad 0 \quad (\text{Falso}, \bot)
\]
L'insieme dei valori di verità è indicato con $\mathbb{V} = \{0, 1\}$.
\end{definizione}

Il calcolo proposizionale classico si fonda su due principi cardine della logica aristotelica:
\begin{enumerate}
    \item \textbf{Principio di Non Contraddizione:} una proposizione non può essere contemporaneamente vera e falsa.
    \item \textbf{Principio del Terzo Escluso (\emph{Tertium non datur}):} una proposizione è necessariamente o vera o falsa; non esiste un terzo valore di verità intermedio.
\end{enumerate}

\section{Connettivi Proposizionali}
I connettivi proposizionali consentono di costruire nuove proposizioni (dette \emph{composte}) a partire da proposizioni date (dette \emph{atomiche}). Formalmente, un connettivo a $n$ argomenti è una funzione $\mathbb{V}^n \to \mathbb{V}$.

\subsection{Connettivo Unario: La Negazione}
L'unico connettivo unario non banale è la \textbf{negazione} ($\neg$ oppure $\sim$):
\[
    \neg: \mathbb{V} \to \mathbb{V}, \quad \neg 1 = 0, \quad \neg 0 = 1
\]

\subsection{Connettivi Binari Fondamentali}
Dati due enunciati $p$ e $q$, i principali connettivi binari sono:
\begin{enumerate}
    \item \textbf{Congiunzione ($\wedge$):} $p \wedge q$ è vera se e solo se sia $p$ che $q$ sono vere.
    \item \textbf{Disgiunzione Inclusiva ($\vee$):} $p \vee q$ è vera se almeno una tra $p$ e $q$ è vera (falsa solo se entrambe sono false).
    \item \textbf{Disgiunzione Esclusiva ($\xor$ o $\dot{\vee}$):} $p \xor q$ è vera se esattamente una tra $p$ e $q$ è vera (corrisponde all'operazione di addizione modulo 2).
    \item \textbf{Implicazione Materiale ($\implica$):} la proposizione condizionale $p \implica q$ (``se $p$, allora $q$'') è definita come equivalente a:
    \[
        (p \implica q) \iff (\neg p \vee q)
    \]
    Di conseguenza, $p \implica q$ è \textbf{falsa unicamente quando l'antecedente $p$ è vero e il conseguente $q$ è falso}. Se la premessa $p$ è falsa, l'implicazione $p \implica q$ è sempre formalmente vera (\emph{ex falso quodlibet}).
    \item \textbf{Coimplicazione o Doppia Implicazione ($\shiff$):} $p \shiff q$ (``$p$ se e solo se $q$'') è vera se $p$ e $q$ possiedono il medesimo valore di verità.
    \item \textbf{Incompatibilità (NAND, $\uparrow$):} $\neg (p \wedge q)$.
    \item \textbf{Negazione congiunta (NOR, $\downarrow$):} $\neg (p \vee q)$.
\end{enumerate}

\begin{table}[htbp]
\centering
\caption{Tavola di verità dei principali connettivi logici binari}
\label{tab:tavola_verita}
\begin{tabular}{|cc||c|c|c|c|c|c|c|}
\hline
$p$ & $q$ & $\neg p$ & $p \wedge q$ & $p \vee q$ & $p \xor q$ & $p \implica q$ & $p \shiff q$ & $p \uparrow q$ \\
\hline
$1$ & $1$ & $0$ & $1$ & $1$ & $0$ & $1$ & $1$ & $0$ \\
$1$ & $0$ & $0$ & $0$ & $1$ & $1$ & $0$ & $0$ & $1$ \\
$0$ & $1$ & $1$ & $0$ & $1$ & $1$ & $1$ & $0$ & $1$ \\
$0$ & $0$ & $1$ & $0$ & $0$ & $0$ & $1$ & $1$ & $1$ \\
\hline
\end{tabular}
\end{table}

\begin{avvisoCutolo}
    \textbf{Attenzione alla definizione di Implicazione Materiale!}
    All'orale il Prof. Cutolo sottolinea costantemente che l'implicazione logica non presuppone alcun legame di causa-effetto tra $p$ e $q$. L'affermazione ``Se la Luna è di formaggio, allora $2+2=5$'' è \textbf{logicamente vera}, poiché la premessa è falsa.
\end{avvisoCutolo}

\section{Tautologie, Contraddizioni ed Equivalenze Logiche}

\begin{definizione}{Tautologia e Contraddizione}{tautologia}
Una formula proposizionale si dice:
\begin{itemize}
    \item \textbf{Tautologia} se assume valore di verità $1$ per qualunque assegnazione dei valori di verità delle variabili proposizionali costituenti.
    \item \textbf{Contraddizione} se assume valore di verità $0$ per qualunque assegnazione.
\end{itemize}
Due formule $P$ e $Q$ si dicono \textbf{logicamente equivalenti} (scritto $P \equiv Q$) se la formula $P \shiff Q$ è una tautologia.
\end{definizione}

\begin{teorema}{Equivalenze Logiche Fondamentali}{equiv_fondamentali}
Per ogni variabile proposizionale $p, q, r$, sussistono le seguenti equivalenze tautologiche:
\begin{enumerate}
    \item \textbf{Idempotenza:} $p \wedge p \equiv p$, \quad $p \vee p \equiv p$.
    \item \textbf{Commutatività:} $p \wedge q \equiv q \wedge p$, \quad $p \vee q \equiv q \vee p$.
    \item \textbf{Associatività:} 
    \[
        (p \wedge q) \wedge r \equiv p \wedge (q \wedge r), \qquad (p \vee q) \vee r \equiv p \vee (q \vee r)
    \]
    \item \textbf{Distributività:}
    \[
        p \wedge (q \vee r) \equiv (p \wedge q) \vee (p \wedge r), \qquad p \vee (q \wedge r) \equiv (p \vee q) \wedge (p \vee r)
    \]
    \item \textbf{Leggi di De Morgan:}
    \[
        \neg(p \wedge q) \equiv \neg p \vee \neg q, \qquad \neg(p \vee q) \equiv \neg p \wedge \neg q
    \]
    \item \textbf{Assorbimento:}
    \[
        p \wedge (p \vee q) \equiv p, \qquad p \vee (p \wedge q) \equiv p
    \]
    \item \textbf{Doppia Negazione:} $\neg(\neg p) \equiv p$.
\end{enumerate}
\end{teorema}

\subsection{Tautologie Notevoli sull'Implicazione}
L'implicazione materiale gode di proprietà cruciali per le dimostrazioni matematiche:
\begin{itemize}
    \item \textbf{Legge di Contrapposizione:}
    \[
        (p \implica q) \equiv (\neg q \implica \neg p)
    \]
    Dimostrare un teorema del tipo ``Se $p$, allora $q$'' equivale logicamente a dimostrare ``Se non $q$, allora non $p$''.
    \item \textbf{Transitività dell'Implicazione (Sillogismo Ipotetico):}
    \[
        ((p \implica q) \wedge (q \implica r)) \implica (p \implica r)
    \]
    \item \textbf{Modus Ponens:} $((p \implica q) \wedge p) \implica q$.
    \item \textbf{Modus Tollens:} $((p \implica q) \wedge \neg q) \implica \neg p$.
    \item \textbf{Legge di Peirce:} $((p \implica q) \implica p) \implica p$ (tautologia valida solo nella logica classica bivalente).
    \item \textbf{Riduzione all'Assurdo:} $(p \implica (q \wedge \neg q)) \implica \neg p$.
\end{itemize}

\section{Calcolo dei Predicati e Quantificatori}

\begin{definizione}{Predicato e Dominio}{predicato}
Un \textbf{predicato} $P(x)$ su un dominio (o universo del discorso) $\mathcal{U}$ è una frase dichiarativa contenente una o più variabili $x \in \mathcal{U}$ che diviene una proposizione non appena a ciascuna variabile viene sostituito un elemento specifico di $\mathcal{U}$.
\end{definizione}

\subsection{I Quantificatori Logici}
Per trasformare un predicato $P(x)$ in una proposizione possiamo ricorrere ai \textbf{quantificatori}:
\begin{enumerate}
    \item \textbf{Quantificatore Universale ($\forall$, ``per ogni''):}
    $\forall x \in \mathcal{U}, P(x)$ asserisce che la proprietà $P(x)$ sussiste per \textbf{tutti} gli elementi del dominio $\mathcal{U}$. È falsa non appena esiste anche un solo elemento $x_0 \in \mathcal{U}$ per cui $P(x_0)$ è falsa (\emph{controesempio}).
    \item \textbf{Quantificatore Esistenziale ($\exists$, ``esiste almeno un''):}
    $\exists x \in \mathcal{U} : P(x)$ asserisce che esiste \textbf{almeno un} elemento in $\mathcal{U}$ che verifica $P(x)$. È vera se si esibisce almeno un elemento soddisfacente.
    \item \textbf{Esistenza ed Unicità ($\exists!$, ``esiste un unico''):}
    $\exists! x \in \mathcal{U} : P(x)$ significa:
    \[
        (\exists x \in \mathcal{U} : P(x)) \wedge (\forall x, y \in \mathcal{U}, (P(x) \wedge P(y)) \implica x = y)
    \]
\end{enumerate}

\subsection{Regole di Negazione dei Quantificatori}
La negazione di un'asserzione quantificata scambia la natura del quantificatore e nega il predicato:
\begin{align}
    \neg (\forall x \in \mathcal{U}, P(x)) &\equiv \exists x \in \mathcal{U} : \neg P(x) \label{eq:neg_all} \\
    \neg (\exists x \in \mathcal{U} : P(x)) &\equiv \forall x \in \mathcal{U}, \neg P(x) \label{eq:neg_ex}
\end{align}

\begin{avvisoCutolo}
    \textbf{L'Ordine dei Quantificatori è Fondamentale!}
    Quantificatori dello stesso tipo commutano: $\forall x \forall y \equiv \forall y \forall x$ e $\exists x \exists y \equiv \exists y \exists x$.
    
    Al contrario, \textbf{quantificatori di tipo diverso non commutano in generale}:
    \[
        \exists y \in \mathcal{U} : \forall x \in \mathcal{U}, P(x, y) \quad \Implica \quad \forall x \in \mathcal{U}, \exists y \in \mathcal{U} : P(x, y)
    \]
    ma l'implicazione inversa è \textbf{falsa}.
    
    \emph{Esempio Classico di Cutolo:} Si consideri il dominio $\N$.
    \begin{itemize}
        \item L'enunciato $\forall x \in \N, \exists y \in \N : y > x$ è \textbf{vero} (per ogni numero naturale esiste un naturale strettamente maggiore, ad es. $y = x+1$).
        \item L'enunciato $\exists y \in \N : \forall x \in \N, y > x$ è \textbf{falso} (asserisce l'esistenza di un singolo numero naturale $y$ maggiore di tutti i numeri naturali, compreso se stesso!).
    \end{itemize}
\end{avvisoCutolo}

\section{Metodologie di Dimostrazione Matematica}
Nel corso di Algebra si incontrano prevalentemente le seguenti strategie di dimostrazione:
\begin{enumerate}
    \item \textbf{Dimostrazione Diretta:} Si assume vera la premessa $P$ e, attraverso una catena di implicazioni logiche giustificate da assiomi o teoremi noti, si giunge a dedurre la tesi $Q$.
    \item \textbf{Dimostrazione per Contrapposizione:} Si sfrutta l'equivalenza $(P \implica Q) \equiv (\neg Q \implica \neg P)$. Si assume $\neg Q$ e si dimostra $\neg P$.
    \item \textbf{Dimostrazione per Assurdo (\emph{Reductio ad absurdum}):} Si assume vera l'ipotesi $P$ e contemporaneamente si nega la tesi $\neg Q$. Da $(P \wedge \neg Q)$ si deriva una contraddizione logica (del tipo $R \wedge \neg R$). Ciò prova che l'assunzione $\neg Q$ era insostenibile, dunque $Q$ è vera.
    \item \textbf{Dimostrazione per Induzione:} Utilizzata per proprietà quantificate su $\N$; la analizzeremo diffusamente nel prossimo capitolo.
    \item \textbf{Dimostrazione per Controesempio:} Per confutare una proposizione universale $\forall x, P(x)$, è sufficiente esibire un singolo elemento $x_0$ tale che $\neg P(x_0)$.
\end{enumerate}

\chapter{Teoria degli Insiemi e Calcolo Combinatorio}
\label{chap:insiemi_combinatoria}

In questo capitolo introduciamo i concetti fondamentali della teoria ingenua degli insiemi, il principio del buon ordinamento di $\N$, il principio di induzione matematica e gli strumenti essenziali del calcolo combinatorio.

\section{Insiemi, Appartenenza e Inclusione}

La teoria degli insiemi assume come nozioni primitive quelle di \textbf{insieme} e di \textbf{appartenenza} ($\in$). Se $x$ è un elemento e $A$ è un insieme, scriviamo $x \in A$ per indicare che $x$ appartiene ad $A$, e $x \notin A$ in caso contrario.

\begin{definizione}{Principio di Estensionalità e Inclusione}{estensionalita}
Due insiemi $A$ e $B$ sono \textbf{uguali} ($A = B$) se e solo se contengono esattamente gli stessi elementi:
\[
    A = B \quad \Shiff \quad \forall x, (x \in A \shiff x \in B)
\]
Dati due insiemi $A$ e $B$, diciamo che $A$ è un \textbf{sottoinsieme} di $B$ (o che $A$ è \textbf{incluso} in $B$), e scriviamo $A \sseq B$, se:
\[
    A \sseq B \quad \Shiff \quad \forall x, (x \in A \implica x \in B)
\]
Si dice che $A$ è un sottoinsieme \textbf{proprio} di $B$ ($A \subset B$) se $A \sseq B$ e $A \neq B$.
\end{definizione}

\begin{proposizione}{Doppia Inclusione}{doppia_inclusione}
Per ogni coppia di insiemi $A$ e $B$, si ha:
\[
    A = B \quad \Shiff \quad (A \sseq B \wedge B \sseq A)
\]
\end{proposizione}

\subsection{L'Insieme Vuoto e lo Schema di Separazione}
Esiste un unico insieme privo di elementi, detto \textbf{insieme vuoto} e denotato con $\vuoto$ oppure $\emptyset$.
Per ogni insieme $A$, vale sempre $\vuoto \sseq A$. Infatti, l'implicazione $x \in \vuoto \implica x \in A$ ha l'antecedente sempre falso, dunque è banalmente vera.

Per evitare contraddizioni logiche (come il celebre \emph{Paradosso di Russell} sull'insieme di tutti gli insiemi che non appartengono a se stessi, $R = \{x \mid x \notin x\}$), non è lecito formare insiemi tramite proprietà arbitrarie senza specificare un insieme universo di riferimento. Si adotta lo \textbf{Schema di Assiomi di Separazione (di Zermelo)}: dato un insieme $A$ e un predicato $P(x)$, è possibile costruire il sottoinsieme:
\[
    \{x \in A \mid P(x)\}
\]

\section{Operazioni Insiemistiche Fondamentali}

Dati due insiemi $A$ e $B$ appartenenti a un prefissato universo $\mathcal{U}$, si definiscono le seguenti operazioni:
\begin{enumerate}
    \item \textbf{Unione:} $A \cup B = \{x \in \mathcal{U} \mid x \in A \vee x \in B\}$.
    \item \textbf{Intersezione:} $A \cap B = \{x \in \mathcal{U} \mid x \in A \wedge x \in B\}$. Se $A \cap B = \vuoto$, $A$ e $B$ si dicono \textbf{disgiunti}.
    \item \textbf{Differenza Insiemistica:} $A \smin B = \{x \in \mathcal{U} \mid x \in A \wedge x \notin B\}$.
    \item \textbf{Complementare (Assoluto):} $\mathcal{C}_\mathcal{U}(A) = A^c = \mathcal{U} \smin A$.
    \item \textbf{Differenza Simmetrica:}
    \[
        A \ds B = (A \smin B) \cup (B \smin A) = (A \cup B) \smin (A \cap B)
    \]
    La differenza simmetrica corrisponde insiemisticamente al connettivo logico $\xor$:
    \[
        x \in A \ds B \quad \Shiff \quad (x \in A) \xor (x \in B)
    \]
\end{enumerate}

\begin{teorema}{Proprietà delle Operazioni Insiemistiche}{prop_insiemi}
Le operazioni insiemistiche godono delle seguenti proprietà:
\begin{enumerate}
    \item \textbf{Leggi di De Morgan:} $(A \cup B)^c = A^c \cap B^c$ \quad e \quad $(A \cap B)^c = A^c \cup B^c$.
    \item \textbf{Distributività:} 
    \[
        A \cap (B \cup C) = (A \cap B) \cup (A \cap C), \qquad A \cup (B \cap C) = (A \cup B) \cap (A \cup C)
    \]
    \item \textbf{Proprietà della Differenza Simmetrica:} $(A \ds B) \ds C = A \ds (B \ds C)$ (associatività), $A \ds B = B \ds A$ (commutatività), $A \ds \vuoto = A$ (elemento neutro), $A \ds A = \vuoto$ (ogni elemento è il simmetrico di se stesso).
\end{enumerate}
\end{teorema}

\section{L'Insieme delle Parti e il Prodotto Cartesiano}

\begin{definizione}{Insieme delle Parti (Power Set)}{power_set}
Dato un insieme $S$, l'\textbf{insieme delle parti} (o booleano di $S$), denotato con $\P(S)$ (oppure $2^S$ o $\mathcal{P}(S)$), è l'insieme di tutti i possibili sottoinsiemi di $S$:
\[
    \P(S) = \{X \mid X \sseq S\}
\]
\end{definizione}

Per qualunque insieme $S$, si ha sempre $\vuoto \in \P(S)$ e $S \in \P(S)$.
\begin{itemize}
    \item Se $S = \vuoto$, allora $\P(\vuoto) = \{\vuoto\}$, dunque $|\P(\vuoto)| = 1$.
    \item Se $S = \{a\}$, allora $\P(S) = \{\vuoto, \{a\}\}$, dunque $|\P(S)| = 2$.
    \item Se $S$ è un insieme finito con $|S| = n$ elementi, allora $|\P(S)| = 2^n$.
\end{itemize}

\subsection{Coppie Ordinate e Prodotto Cartesiano}
Nella teoria degli insiemi l'ordine degli elementi non conta: $\{a, b\} = \{b, a\}$. Per definire una coppia in cui l'ordine sia rilevante, si usa la \textbf{definizione di Kuratowski}:
\[
    (a, b) \coloneqq \{\{a\}, \{a, b\}\}
\]
Tale definizione assicura la proprietà caratteristica fondamentale:
\[
    (a, b) = (c, d) \quad \Shiff \quad (a = c \wedge b = d)
\]

\begin{definizione}{Prodotto Cartesiano}{prod_cartesiano}
Dati due insiemi $A$ e $B$, il loro \textbf{prodotto cartesiano} $A \times B$ è l'insieme delle coppie ordinate:
\[
    A \times B = \{(a, b) \mid a \in A \wedge b \in B\}
\]
Se $A$ o $B$ è vuoto, allora $A \times B = \vuoto$. Se $A$ e $B$ sono finiti, $|A \times B| = |A| \cdot |B|$.
\end{definizione}

\section{Il Buon Ordinamento di $\N$ e il Principio di Induzione}

Ricordiamo la convenzione adottata nel corso del Prof. Cutolo: $\N = \{0, 1, 2, 3, \dots\}$.

\begin{teorema}{Principio del Buon Ordinamento di $\N$}{buon_ord}
Ogni sottoinsieme non vuoto $S \sseq \N$ ammette \textbf{minimo}, ovvero esiste un elemento $m \in S$ tale che $m \le s$ per ogni $s \in S$.
\end{teorema}

Dal buon ordinamento deriva la tecnica di dimostrazione per \textbf{controesempio minimo}: per dimostrare che una proprietà $P(n)$ vale per ogni $n \in \N$, si suppone per assurdo che l'insieme dei controesempi $C = \{n \in \N \mid \neg P(n)\}$ sia non vuoto. Per il principio del buon ordinamento, $C$ deve ammettere un minimo elemento $m = \min(C)$. Analizzando le proprietà di $m$, si deduce una contraddizione logica.

\begin{teorema}{Principio di Induzione Matematica}{induzione}
Sia $P(n)$ un predicato definito per $n \in \N$. Supponiamo che:
\begin{enumerate}
    \item \textbf{Base dell'induzione:} $P(0)$ è vera;
    \item \textbf{Passo induttivo:} per ogni $n \in \N$, se $P(n)$ è vera, allora $P(n+1)$ è vera:
    \[
        \forall n \in \N, (P(n) \implica P(n+1))
    \]
\end{enumerate}
Allora $P(n)$ è vera per ogni $n \in \N$.
\end{teorema}

\begin{proof}
Supponiamo per assurdo che non per tutti gli $n \in \N$ valga $P(n)$. Allora l'insieme dei controesempi $S = \{k \in \N \mid \neg P(k)\}$ è non vuoto. Per il Principio del Buon Ordinamento, $S$ ammette un minimo elemento $m \in S$.
Poiché per ipotesi (1) $P(0)$ è vera, deduciamo che $0 \notin S$, dunque $m > 0$. Ne segue che $m - 1 \in \N$.
Essendo $m - 1 < m$, ed essendo $m$ il minimo di $S$, deve necessariamente accadere che $m - 1 \notin S$, vale a dire che $P(m-1)$ è vera.
Ma per l'ipotesi induttiva (2), $P(m-1) \implica P((m-1)+1) = P(m)$. Dunque $P(m)$ è vera, il che contraddice l'assunzione che $m \in S$.
L'assurdo prova che $S = \vuoto$, quindi $P(n)$ è vera per ogni $n \in \N$.
\end{proof}

\subsection{Induzione Completa (o Forte)}
Nella forma forte dell'induzione, il passo induttivo assume come ipotesi che la proprietà sia vera per \textbf{tutti} gli interi $k \le n$:
\[
    (\forall k \le n, P(k)) \Implica P(n+1)
\]
Le due forme sono logicamente equivalenti tra loro e al principio del buon ordinamento.

\section{Elementi di Calcolo Combinatorio}

Sia $S$ un insieme finito. Denotiamo con $|S|$ la sua \textbf{cardinalità} (il numero dei suoi elementi).

\subsection{Principi Fondamentali del Conteggio}
\begin{itemize}
    \item \textbf{Regola della Somma:} se $A$ e $B$ sono insiemi finiti disgiunti ($A \cap B = \vuoto$), allora $|A \cup B| = |A| + |B|$.
    \item \textbf{Regola del Prodotto:} per due insiemi finiti qualsiasi, $|A \times B| = |A| \cdot |B|$.
    \item \textbf{Principio di Inclusione-Esclusione (PIE):} per due insiemi qualsiasi:
    \[
        |A \cup B| = |A| + |B| - |A \cap B|
    \]
    Per tre insiemi $A, B, C$:
    \[
        |A \cup B \cup C| = |A| + |B| + |C| - |A \cap B| - |A \cap C| - |B \cap C| + |A \cap B \cap C|
    \]
\end{itemize}

\subsection{Disposizioni, Permutazioni e Combinazioni}
Sia $S$ un insieme di $n$ elementi. Consideriamo il problema di selezionare $k$ elementi da $S$:

\begin{enumerate}
    \item \textbf{Disposizioni Semplici ($D_{n, k}$):} sequenze ordinate di $k$ elementi \emph{distinti} scelti tra gli $n$ disponibili ($0 \le k \le n$):
    \[
        D_{n, k} = n(n-1)(n-2)\cdots(n-k+1) = \frac{n!}{(n-k)!} = n^{\underline{k}}
    \]
    \item \textbf{Permutazioni Semplici ($P_n$):} disposizioni di tutti gli $n$ elementi ($k = n$):
    \[
        P_n = D_{n, n} = n! = n \cdot (n-1) \cdots 2 \cdot 1, \qquad \text{con } 0! = 1
    \]
    \item \textbf{Disposizioni con Ripetizione ($D'_{n, k}$):} sequenze ordinate di $k$ elementi in cui ciascun elemento può ripetersi:
    \[
        D'_{n, k} = n^k = |A^B| \quad (\text{con } |A| = n, |B| = k)
    \]
    Esso conta il numero totale di applicazioni da un insieme di cardinalità $k$ a un insieme di cardinalità $n$.
    \item \textbf{Combinazioni Semplici ($C_{n, k}$ o $\binom{n}{k}$):} sottoinsiemi non ordinati di $k$ elementi distinti scelti tra gli $n$ disponibili:
    \[
        \binom{n}{k} = \frac{D_{n, k}}{P_k} = \frac{n!}{k!(n-k)!}
    \]
    Il simbolo $\binom{n}{k}$ si legge ``$n$ su $k$'' ed è detto \textbf{coefficiente binomiale}.
\end{enumerate}

\subsection{Proprietà dei Coefficienti Binomiali e Triangolo di Tartaglia}

\begin{teorema}{Proprietà dei Coefficienti Binomiali}{prop_binomio}
Per ogni $n, k \in \N$ con $0 \le k \le n$:
\begin{enumerate}
    \item \textbf{Simmetria:} $\binom{n}{k} = \binom{n}{n-k}$.
    \item \textbf{Formula di Ricorrenza di Tartaglia-Pascal:}
    \[
        \binom{n}{k} = \binom{n-1}{k-1} + \binom{n-1}{k} \quad (\text{per } 1 \le k \le n-1)
    \]
    \item \textbf{Somma per riga (Numero totale di sottoinsiemi):}
    \[
        \sum_{k=0}^n \binom{n}{k} = \binom{n}{0} + \binom{n}{1} + \dots + \binom{n}{n} = 2^n = |\P(S)|
    \]
\end{enumerate}
\end{teorema}

\begin{proof}[Dimostrazione combinatoria della formula di Tartaglia-Pascal]
Sia $S$ un insieme di $n$ elementi e fissiamo un elemento arbitrario $x_0 \in S$. Vogliamo contare i sottoinsiemi di $S$ aventi esattamente $k$ elementi (il cui numero è per definizione $\binom{n}{k}$).
I $k$-sottoinsiemi di $S$ si partizionano naturalmente in due classi disgiunte:
\begin{itemize}
    \item I sottoinsiemi che \textbf{contengono} $x_0$: per formarli, dobbiamo scegliere i restanti $k-1$ elementi tra gli $n-1$ elementi di $S \smin \{x_0\}$. Ce ne sono $\binom{n-1}{k-1}$.
    \item I sottoinsiemi che \textbf{non contengono} $x_0$: sono formati scegliendo tutti i $k$ elementi direttamente in $S \smin \{x_0\}$. Ce ne sono $\binom{n-1}{k}$.
\end{itemize}
Per la regola della somma, $\binom{n}{k} = \binom{n-1}{k-1} + \binom{n-1}{k}$.
\end{proof}

\begin{teorema}{Formula del Binomio di Newton}{binomio_newton}
Dati due elementi $a, b$ appartenenti a un anello commutativo (o numeri reali), per ogni $n \in \N$ vale:
\[
    (a + b)^n = \sum_{k=0}^n \binom{n}{k} a^{n-k} b^k = \binom{n}{0}a^n + \binom{n}{1}a^{n-1}b + \dots + \binom{n}{n}b^n
\]
\end{teorema}

\subsection{Il Principio dei Cassetti (Pigeonhole Principle)}
\begin{teorema}{Principio dei Cassetti}{pigeonhole}
Se si collocano $n$ oggetti in $k$ contenitori (cassetti) con $n > k$, allora almeno un cassetto conterrà almeno $\lceil n/k \rceil \ge 2$ oggetti.
\end{teorema}
Questo principio, apparentemente elementare, è uno strumento dimostrativo potentissimo in algebra (ad esempio per provare che ogni applicazione iniettiva tra insiemi finiti della stessa cardinalità è suriettiva, o che ogni dominio d'integrità finito è un campo).


% ------------------------------------------------------------------------------
% PARTE II: RELAZIONI, APPLICAZIONI E RETICOLI
% ------------------------------------------------------------------------------
\part{Relazioni, Applicazioni, Ordini e Strutture Booleane}
\chapter{Relazioni Binarie, Applicazioni e Cardinalità}
\label{chap:relazioni_applicazioni}

In questo capitolo affrontiamo uno dei cardini dell'intero programma di Algebra: la teoria delle corrispondenze e delle relazioni binarie, lo studio formale delle applicazioni, la teoria delle sezioni e retrazioni, il Teorema fondamentale di omomorfismo per insiemi e la nozione di cardinalità.

\section{Corrispondenze e Relazioni Binarie}

\begin{definizione}{Corrispondenza}{corrispondenza}
Dati due insiemi $A$ e $B$, una \textbf{corrispondenza} da $A$ a $B$ è una terna ordinata $\mathcal{C} = (A, B, G)$, dove:
\begin{itemize}
    \item $A$ è detto \textbf{insieme di partenza} (o dominio);
    \item $B$ è detto \textbf{insieme di arrivo} (o codominio);
    \item $G \sseq A \times B$ è un sottoinsieme del prodotto cartesiano, detto \textbf{grafo} della corrispondenza.
\end{itemize}
Se $(a, b) \in G$, si scrive frequentemente $a \mathcal{C} b$ e si dice che $a$ è in corrispondenza con $b$.

La \textbf{corrispondenza inversa} di $\mathcal{C}$ è la terna $\mathcal{C}^{-1} = (B, A, G^{-1})$, ove:
\[
    G^{-1} = \{(b, a) \in B \times A \mid (a, b) \in G\}
\]
\end{definizione}

\begin{definizione}{Relazione Binaria}{relazione_binaria}
Una corrispondenza in cui l'insieme di partenza coincide con l'insieme di arrivo, cioè del tipo $\mathcal{R} = (A, A, G)$ con $G \sseq A \times A$, è detta \textbf{relazione binaria omogenea} sull'insieme $A$.
L'insieme di tutte le relazioni binarie su $A$ si denota con $\Rel(A) = \P(A \times A)$. Se $|A| = n$, allora $|\Rel(A)| = 2^{n^2}$.
\end{definizione}

\subsection{Proprietà delle Relazioni Binarie}
Sia $\mathcal{R}$ una relazione binaria su un insieme $A$. Essa può godere delle seguenti proprietà:
\begin{enumerate}
    \item \textbf{Riflessività:} ogni elemento è in relazione con se stesso:
    \[
        \forall x \in A, \quad x \mathcal{R} x \quad (\text{ovvero } \id_A \sseq G)
    \]
    \item \textbf{Antiriflessività:} nessun elemento è in relazione con se stesso:
    \[
        \forall x \in A, \quad \neg(x \mathcal{R} x) \quad (\text{ovvero } \id_A \cap G = \vuoto)
    \]
    \begin{avvisoCutolo}
        \textbf{Antiriflessivo NON significa ``non riflessivo''!}
        Una relazione non riflessiva è una relazione per cui \emph{esiste almeno un} elemento $x \in A$ tale che $\neg(x \mathcal{R} x)$. Una relazione antiriflessiva richiede invece che \emph{per tutti} gli elementi valga $\neg(x \mathcal{R} x)$. Esistono relazioni che non sono né riflessive né antiriflessive (ad esempio quando alcuni elementi sono in relazione con se stessi ed altri no).
    \end{avvisoCutolo}
    \item \textbf{Simmetria:} se $x$ è in relazione con $y$, allora anche $y$ è in relazione con $x$:
    \[
        \forall x, y \in A, \quad x \mathcal{R} y \Implica y \mathcal{R} x \quad (\text{ovvero } G = G^{-1})
    \]
    \item \textbf{Antisimmetria:} due elementi distinti non possono essere reciprocamente in relazione:
    \[
        \forall x, y \in A, \quad (x \mathcal{R} y \wedge y \mathcal{R} x) \Implica x = y \quad (\text{ovvero } G \cap G^{-1} \sseq \id_A)
    \]
    Equivalentemente, per contrapposizione: $\forall x \neq y, x \mathcal{R} y \implica \neg(y \mathcal{R} x)$.
    \item \textbf{Transitività:} se $x$ è in relazione con $y$ e $y$ è in relazione con $z$, allora $x$ è in relazione con $z$:
    \[
        \forall x, y, z \in A, \quad (x \mathcal{R} y \wedge y \mathcal{R} z) \Implica x \mathcal{R} z \quad (\text{ovvero } G \circ G \sseq G)
    \]
    \item \textbf{Totalità (o Connessione):} per ogni coppia di elementi, almeno uno è in relazione con l'altro:
    \[
        \forall x, y \in A, \quad x \mathcal{R} y \vee y \mathcal{R} x
    \]
\end{enumerate}

\subsection{Prodotto Relazionale}
Date due corrispondenze $\mathcal{R} = (A, B, G_\mathcal{R})$ e $\mathcal{S} = (B, C, G_\mathcal{S})$, il loro \textbf{prodotto relazionale} (o composizione) $\mathcal{S} \circ \mathcal{R}$ è la corrispondenza $(A, C, G_{\mathcal{S}\circ\mathcal{R}})$ definita da:
\[
    x (\mathcal{S} \circ \mathcal{R}) z \quad \Shiff \quad \exists y \in B : (x \mathcal{R} y \wedge y \mathcal{S} z)
\]
Il prodotto relazionale è un'operazione \textbf{associativa}: per ogni terna di relazioni componibili vale $\mathcal{T} \circ (\mathcal{S} \circ \mathcal{R}) = (\mathcal{T} \circ \mathcal{S}) \circ \mathcal{R}$.

\section{Relazioni di Equivalenza e Partizioni}

\begin{definizione}{Relazione di Equivalenza}{rel_equivalenza}
Una relazione binaria $\mathcal{R}$ su un insieme $A$ si dice \textbf{relazione di equivalenza} se è:
\begin{enumerate}
    \item \textbf{Riflessiva:} $\forall x \in A, x \mathcal{R} x$;
    \item \textbf{Simmetrica:} $\forall x, y \in A, x \mathcal{R} y \implica y \mathcal{R} x$;
    \item \textbf{Transitiva:} $\forall x, y, z \in A, (x \mathcal{R} y \wedge y \mathcal{R} z) \implica x \mathcal{R} z$.
\end{enumerate}
L'insieme di tutte le relazioni di equivalenza su $A$ si denota con $\Eq(A)$.
\end{definizione}

\begin{definizione}{Classe di Equivalenza e Insieme Quoziente}{classe_eq}
Sia $\mathcal{R} \in \Eq(A)$ e sia $a \in A$. La \textbf{classe di equivalenza} determinata da $a$ modulo $\mathcal{R}$ è il sottoinsieme degli elementi di $A$ equivalenti ad $a$:
\[
    [a]_\mathcal{R} = \{x \in A \mid x \mathcal{R} a\}
\]
L'elemento $a$ è detto \textbf{rappresentante} della classe $[a]_\mathcal{R}$.
L'insieme avente come elementi tutte le classi di equivalenza è detto \textbf{insieme quoziente} di $A$ rispetto a $\mathcal{R}$ e si denota con:
\[
    A/\mathcal{R} = \{[a]_\mathcal{R} \mid a \in A\}
\]
L'applicazione suriettiva $\pi_\mathcal{R}: A \to A/\mathcal{R}$ definita da $\pi_\mathcal{R}(a) = [a]_\mathcal{R}$ è detta \textbf{proiezione canonica}.
\end{definizione}

\begin{proposizione}{Proprietà delle Classi di Equivalenza}{prop_classi}
Sia $\mathcal{R}$ una relazione di equivalenza su $A$. Allora:
\begin{enumerate}
    \item Per ogni $a \in A$, si ha $a \in [a]_\mathcal{R}$ (ogni classe è non vuota).
    \item Per ogni coppia $a, b \in A$, le seguenti affermazioni sono equivalenti:
    \[
        a \mathcal{R} b \quad \Shiff \quad [a]_\mathcal{R} = [b]_\mathcal{R} \quad \Shiff \quad [a]_\mathcal{R} \cap [b]_\mathcal{R} \neq \vuoto
    \]
    \item Due classi di equivalenza o coincidono oppure sono completamente disgiunte:
    \[
        [a]_\mathcal{R} \neq [b]_\mathcal{R} \quad \Implica \quad [a]_\mathcal{R} \cap [b]_\mathcal{R} = \vuoto
    \]
    \item L'unione di tutte le classi di equivalenza coincide con l'intero insieme $A$:
    \[
        \bigcup_{a \in A} [a]_\mathcal{R} = A
    \]
\end{enumerate}
\end{proposizione}

\subsection{Partizioni di un Insieme}

\begin{definizione}{Partizione}{partizione}
Una \textbf{partizione} di un insieme non vuoto $A$ è una famiglia $\mathcal{F} \sseq \P(A)$ di sottoinsiemi di $A$ (detti \emph{blocchi} della partizione) tale che:
\begin{enumerate}
    \item Nessun blocco è vuoto: $\forall X \in \mathcal{F}, X \neq \vuoto$;
    \item I blocchi sono a due a due disgiunti: $\forall X, Y \in \mathcal{F}, (X \neq Y \implica X \cap Y = \vuoto)$;
    \item L'unione di tutti i blocchi ricopre $A$: $\bigcup_{X \in \mathcal{F}} X = A$.
\end{enumerate}
L'insieme di tutte le partizioni di $A$ si denota con $\Part(A)$.
\end{definizione}

\begin{teorema}{Teorema Fondamentale delle Relazioni di Equivalenza}{teo_fond_equiv}
Per ogni insieme $A$, esiste una \textbf{biiezione naturale} tra l'insieme delle relazioni di equivalenza su $A$ e l'insieme delle partizioni di $A$:
\[
    \Phi: \Eq(A) \longrightarrow \Part(A), \qquad \mathcal{R} \longmapsto A/\mathcal{R}
\]
La cui applicazione inversa $\Psi: \Part(A) \to \Eq(A)$ associa a ciascuna partizione $\mathcal{F}$ la relazione d'equivalenza $\mathcal{R}_\mathcal{F}$ definita da:
\[
    x \mathcal{R}_\mathcal{F} y \quad \Shiff \quad \exists X \in \mathcal{F} : (x \in X \wedge y \in X)
\]
\end{teorema}

\begin{proof}
La dimostrazione si articola in due verifiche:
\begin{enumerate}
    \item \emph{$\Phi$ è ben posta:} Per la Proposizione~\ref{prop:prop_classi}, per ogni $\mathcal{R} \in \Eq(A)$, l'insieme quoziente $A/\mathcal{R}$ è formato da blocchi non vuoti, disgiunti e che ricoprono $A$. Quindi $A/\mathcal{R} \in \Part(A)$.
    \item \emph{$\Psi$ è ben posta:} Sia $\mathcal{F} \in \Part(A)$. La relazione $\mathcal{R}_\mathcal{F}$ è:
    \begin{itemize}
        \item \emph{Riflessiva:} poiché $\bigcup_{X \in \mathcal{F}} X = A$, ogni $x \in A$ appartiene a qualche $X \in \mathcal{F}$, dunque $x$ e $x$ stanno nello stesso blocco, cioè $x \mathcal{R}_\mathcal{F} x$.
        \item \emph{Simmetrica:} banale per commutatività logica di $x \in X \wedge y \in X$.
        \item \emph{Transitiva:} se $x \mathcal{R}_\mathcal{F} y$ e $y \mathcal{R}_\mathcal{F} z$, esistono $X, Y \in \mathcal{F}$ tali che $x, y \in X$ e $y, z \in Y$. Allora $y \in X \cap Y$, dunque $X \cap Y \neq \vuoto$. Essendo $\mathcal{F}$ una partizione, deve essere $X = Y$. Ne segue che $x, z \in X$, quindi $x \mathcal{R}_\mathcal{F} z$.
    \end{itemize}
    \item \emph{Composizione inversa:} Verifichiamo che $\Psi \circ \Phi = \id_{\Eq(A)}$ e $\Phi \circ \Psi = \id_{\Part(A)}$:
    Data $\mathcal{R} \in \Eq(A)$, $x \mathcal{R}_{A/\mathcal{R}} y \iff \exists [a] \in A/\mathcal{R} : x, y \in [a] \iff [x] = [y] \iff x \mathcal{R} y$.
    Data $\mathcal{F} \in \Part(A)$, per ogni $X \in \mathcal{F}$ e $x \in X$, la classe $[x]_{\mathcal{R}_\mathcal{F}}$ coincide esattamente con l'unico blocco $X$ contenente $x$.
\end{enumerate}
\end{proof}

\section{Applicazioni (Funzioni) tra Insiemi}

\begin{definizione}{Applicazione}{applicazione}
Un'applicazione (o funzione) $f: A \to B$ è una corrispondenza $(A, B, G_f)$ che è:
\begin{itemize}
    \item \textbf{Ovunque definita:} per ogni elemento $a \in A$ esiste almeno un $b \in B$ tale che $(a, b) \in G_f$;
    \item \textbf{Univalente (o funzionale):} se $(a, b) \in G_f$ e $(a, c) \in G_f$, allora $b = c$.
\end{itemize}
In sintesi: $\forall a \in A, \exists! b \in B : (a, b) \in G_f$. Tale unico elemento $b$ si indica con $f(a)$ ed è detto \textbf{immagine} di $a$ tramite $f$.
L'insieme di tutte le applicazioni da $A$ a $B$ si indica con $\Map(A, B)$ o $B^A$. Se $|A| = m$ e $|B| = n$, allora $|B^A| = n^m$.
\end{definizione}

\subsection{Immagine e Controimmagine di Sottoinsiemi}
Data $f: A \to B$:
\begin{itemize}
    \item Per ogni $X \sseq A$, l'\textbf{immagine} di $X$ è l'insieme:
    \[
        \vecf(X) = \{f(x) \in B \mid x \in X\}
    \]
    In particolare, $\im(f) = \vecf(A)$ è l'immagine globale di $f$.
    \item Per ogni $Y \sseq B$, l'\textbf{anti-immagine} (o controimmagine) di $Y$ è l'insieme:
    \[
        \antivecf(Y) = \{x \in A \mid f(x) \in Y\}
    \]
\end{itemize}

\begin{teorema}{Comportamento delle Operazioni Insiemistiche sotto $\vecf$ e $\antivecf$}{teo_immagini}
Sia $f: A \to B$. Per ogni $X_1, X_2 \sseq A$ e ogni $Y_1, Y_2 \sseq B$:
\begin{enumerate}
    \item $\vecf(X_1 \cup X_2) = \vecf(X_1) \cup \vecf(X_2)$.
    \item $\vecf(X_1 \cap X_2) \sseq \vecf(X_1) \cap \vecf(X_2)$ (l'uguaglianza vale in generale \textbf{se e solo se $f$ è iniettiva}).
    \item $\antivecf(Y_1 \cup Y_2) = \antivecf(Y_1) \cup \antivecf(Y_2)$.
    \item $\antivecf(Y_1 \cap Y_2) = \antivecf(Y_1) \cap \antivecf(Y_2)$.
    \item $\antivecf(B \smin Y) = A \smin \antivecf(Y)$ (la controimmagine preserva il complemento).
\end{enumerate}
\end{teorema}

\begin{avvisoCutolo}
    La controimmagine $\antivecf: \P(B) \to \P(A)$ è un omomorfismo completo di strutture booleane: commuta sia con l'unione che con l'intersezione e il complemento!
    L'immagine diretta $\vecf$, invece, commuta solo con l'unione: per l'intersezione vale solo l'inclusione $\vecf(X_1 \cap X_2) \sseq \vecf(X_1) \cap \vecf(X_2)$.
\end{avvisoCutolo}

\subsection{Iniettività, Suriettività e Biiettività}
Sia $f: A \to B$:
\begin{enumerate}
    \item $f$ è \textbf{iniettiva} se elementi distinti hanno immagini distinte:
    \[
        \forall x_1, x_2 \in A, \quad f(x_1) = f(x_2) \Implica x_1 = x_2
    \]
    Equivalentemente, $\forall y \in B$, la fibra $|\antivecf(\{y\})| \le 1$.
    \item $f$ è \textbf{suriettiva} se ogni elemento del codominio è immagine di almeno un elemento del dominio:
    \[
        \forall y \in B, \quad \exists x \in A : f(x) = y \quad (\text{cioè } \im(f) = B)
    \]
    Equivalentemente, $\forall y \in B$, la fibra $|\antivecf(\{y\})| \ge 1$.
    \item $f$ è \textbf{biettiva} (o una biiezione) se è sia iniettiva che suriettiva. In tal caso, $\forall y \in B, |\antivecf(\{y\})| = 1$.
\end{enumerate}

\begin{teorema}{Applicazioni tra Insiemi Finiti di Uguale Cardinalità}{teo_finiti_uguali}
Siano $A$ e $B$ insiemi finiti tali che $|A| = |B|$. Per ogni applicazione $f: A \to B$, le seguenti condizioni sono equivalenti:
\[
    f \text{ è iniettiva} \quad \Shiff \quad f \text{ è suriettiva} \quad \Shiff \quad f \text{ è biettiva}
\]
\end{teorema}

\section{Sezioni, Retrazioni e Applicazioni Invertibili}

Siano $f: A \to B$ e $g: B \to C$. L'applicazione composta $g \circ f: A \to C$ è definita da $(g \circ f)(a) = g(f(a))$. La composizione è associativa e ammette le applicazioni identiche $\id_A, \id_B$ come elementi neutri.

\begin{definizione}{Sezioni e Retrazioni}{sezioni_retrazioni}
Sia $f: A \to B$ un'applicazione con $A \neq \vuoto$. Un'applicazione $g: B \to A$ si dice:
\begin{itemize}
    \item Una \textbf{sezione} di $f$ se $f \circ g = \id_B$ ($g$ è inversa a destra di $f$).
    \item Una \textbf{retrazione} di $f$ se $g \circ f = \id_A$ ($g$ è inversa a sinistra di $f$).
\end{itemize}
\end{definizione}

\begin{teorema}{Caratterizzazione di Iniettività e Suriettività tramite Sezioni e Retrazioni}{teo_sez_retraz}
Sia $f: A \to B$ con $A \neq \vuoto$.
\begin{enumerate}
    \item $f$ è \textbf{suriettiva} se e solo se ammette almeno una \textbf{sezione} $s: B \to A$.
    \item $f$ è \textbf{iniettiva} se e solo se ammette almeno una \textbf{retrazione} $r: B \to A$.
    \item $f$ è \textbf{biettiva} se e solo se è \textbf{invertibile}, ovvero ammette un'unica applicazione $f^{-1}: B \to A$ che è contemporaneamente sezione e retrazione:
    \[
        f \circ f^{-1} = \id_B \quad \text{e} \quad f^{-1} \circ f = \id_A
    \]
\end{enumerate}
\end{teorema}

\begin{proof}
Dimostriamo le due equivalenze fondamentali:
\begin{enumerate}
    \item $(\implica)$ Sia $f$ suriettiva. Allora per ogni $b \in B$, la fibra $\antivecf(\{b\}) \neq \vuoto$. Scegliendo per ciascun $b \in B$ un elemento $x_b \in \antivecf(\{b\})$ e ponendo $s(b) \coloneqq x_b$, si ottiene un'applicazione $s: B \to A$ tale che $(f \circ s)(b) = f(s(b)) = f(x_b) = b = \id_B(b)$. Dunque $s$ è una sezione.
    $(\impliedby)$ Se esiste $s: B \to A$ tale che $f \circ s = \id_B$, allora per ogni $b \in B$ si ha $b = f(s(b))$, dunque $b \in \im(f)$ con preimmagine $s(b)$. Quindi $f$ è suriettiva.
    \item $(\implica)$ Sia $f$ iniettiva. Poiché $A \neq \vuoto$, fissiamo un elemento arbitrario $a_0 \in A$. Definiamo $r: B \to A$ ponendo:
    \[
        r(b) = \begin{cases}
            a & \text{se } b \in \im(f), \text{ dove } a \in A \text{ è l'unico elemento tale che } f(a) = b \\
            a_0 & \text{se } b \notin \im(f)
        \end{cases}
    \]
    Per ogni $a \in A$, $(r \circ f)(a) = r(f(a)) = a = \id_A(a)$, quindi $r$ è una retrazione.
    $(\impliedby)$ Se esiste $r$ con $r \circ f = \id_A$, siano $x_1, x_2 \in A$ tali che $f(x_1) = f(x_2)$. Applicando $r$ a entrambi i membri: $r(f(x_1)) = r(f(x_2)) \implies \id_A(x_1) = \id_A(x_2) \implies x_1 = x_2$. Quindi $f$ è iniettiva.
\end{enumerate}
\end{proof}

\section{Teorema di Omomorfismo per Insiemi}

\begin{definizione}{Nucleo di Equivalenza di un'Applicazione}{nucleo_eq}
Data un'applicazione $f: A \to B$, la relazione $\ker(f)$ (denotata anche con $\sim_f$ o $\mathcal{R}_f$) definita su $A$ da:
\[
    x \sim_f y \quad \Shiff \quad f(x) = f(y)
\]
è una relazione di equivalenza su $A$, detta \textbf{nucleo di equivalenza} associato ad $f$.
La classe di equivalenza di $a \in A$ è la fibra $[a]_{\sim_f} = \antivecf(\{f(a)\})$.
\end{definizione}

\begin{teorema}{Teorema Fondamentale di Omomorfismo per Insiemi}{teo_omom_insiemi}
Sia $f: A \to B$ un'applicazione. Esiste un'\textbf{unica biiezione} $\tilde{f}: A/\ker(f) \to \im(f)$ tale che il seguente diagramma commuti:
\[
\begin{tikzcd}[row sep=large, column sep=large]
A \arrow[r, "f"] \arrow[d, "\pi", swap] & B \\
A/\ker(f) \arrow[r, "\tilde{f}", swap] & \im(f) \arrow[u, "j", swap, hook]
\end{tikzcd}
\]
ovvero tale che $f = j \circ \tilde{f} \circ \pi$, dove:
\begin{itemize}
    \item $\pi: A \to A/\ker(f)$ è la proiezione canonica sul quoziente (suriettiva);
    \item $\tilde{f}: [a] \mapsto f(a)$ è la biiezione canonica;
    \item $j: \im(f) \hookrightarrow B$ è l'immersione canonica (iniettiva).
\end{itemize}
\end{teorema}

\begin{proof}
Verifichiamo che $\tilde{f}: A/\ker(f) \to \im(f)$, definita da $\tilde{f}([a]) = f(a)$, è ben posta e biettiva:
\begin{enumerate}
    \item \emph{Buona positura:} siano $a, b \in A$ con $[a] = [b]$. Allora per definizione $a \sim_f b \implies f(a) = f(b)$. Quindi il valore $\tilde{f}([a])$ non dipende dal rappresentante scelto.
    \item \emph{Iniettività:} $\tilde{f}([a]) = \tilde{f}([b]) \implies f(a) = f(b) \implies a \sim_f b \implies [a] = [b]$.
    \item \emph{Suriettività:} per ogni $y \in \im(f)$, esiste $x \in A$ tale che $f(x) = y$. Allora $\tilde{f}([x]) = f(x) = y$.
\end{enumerate}
\end{proof}

\section{Equipotenza e Cardinalità}

\begin{definizione}{Equipotenza}{equipotenza}
Due insiemi $A$ e $B$ si dicono \textbf{equipotenti} (e si scrive $A \sim B$ o $|A| = |B|$) se esiste una biiezione $f: A \to B$.
L'equipotenza è una relazione di equivalenza sulla classe di tutti gli insiemi.
\end{definizione}

Un insieme $A$ si dice:
\begin{itemize}
    \item \textbf{Finito} se $A = \vuoto$ oppure esiste $n \in \N^*$ tale che $A \sim \{1, 2, \dots, n\}$.
    \item \textbf{Infinito} se non è finito (equivalentemente, se è equipotente a un suo sottoinsieme proprio, secondo la definizione di Dedekind).
    \item \textbf{Numerabile} se $A \sim \N$.
    \item \textbf{Al più numerabile} se è finito oppure numerabile.
    \item \textbf{Continuo} se $A \sim \R$.
\end{itemize}

Gli insiemi $\Z$ e $\Q$ sono numerabili ($|\Z| = |\Q| = \aleph_0$). Al contrario, $\R$ ha la potenza del continuo ($2^{\aleph_0} > \aleph_0$) e non è numerabile.

\begin{teorema}{Teorema di Cantor}{cantor}
Per ogni insieme $A$, non esiste alcuna applicazione \textbf{suriettiva} da $A$ a $\P(A)$. Di conseguenza, $|A| < |\P(A)|$.
\end{teorema}

\begin{proof}[Dimostrazione con l'Argomento Diagonale di Cantor]
Sia $f: A \to \P(A)$ un'applicazione arbitraria. Vogliamo dimostrare che $\im(f) \neq \P(A)$.
Definiamo il sottoinsieme diagonale $D \sseq A$ come:
\[
    D = \{x \in A \mid x \notin f(x)\}
\]
Chiaramente $D \in \P(A)$. Mostriamo che $D$ non è immagine di alcun elemento di $A$.
Supponiamo per assurdo che esista $d \in A$ tale che $f(d) = D$.
Poniamoci la domanda se $d \in D$:
\begin{itemize}
    \item Se $d \in D$, allora per definizione di $D$ si deve avere $d \notin f(d)$. Ma $f(d) = D$, quindi $d \notin D$, contraddizione!
    \item Se $d \notin D$, allora per definizione di $D$ deve essere falso che $d \notin f(d)$, dunque $d \in f(d) = D$, contraddizione!
\end{itemize}
In entrambi i casi si giunge a una palese contraddizione logica. Ne segue che $D \notin \im(f)$, dunque $f$ non può essere suriettiva.
\end{proof}

\chapter{Insiemi Ordinati, Reticoli e Strutture Booleane}
\label{chap:ordini_reticoli}

Lo studio delle relazioni d'ordine parziale e dei reticoli costituisce uno dei temi più ricorrenti e qualificanti delle prove scritte e dei colloqui orali con il Prof. Cutolo. In questo capitolo sviluppiamo la teoria a partire dai diagrammi di Hasse fino al teorema di Stone per le algebre di Boole finite.

\section{Insiemi Parzialmente Ordinati (Poset)}

\begin{definizione}{Relazione d'Ordine}{rel_ordine}
Una relazione binaria $\le$ su un insieme $P$ è detta \textbf{relazione d'ordine parziale (largo)} se soddisfa:
\begin{enumerate}
    \item \textbf{Riflessività:} $\forall x \in P, x \le x$;
    \item \textbf{Antisimmetria:} $\forall x, y \in P, (x \le y \wedge y \le x) \implica x = y$;
    \item \textbf{Transitività:} $\forall x, y, z \in P, (x \le y \wedge y \le z) \implica x \le z$.
\end{enumerate}
La coppia $(P, \le)$ è detta \textbf{insieme parzialmente ordinato} (o \textbf{poset}, dall'inglese \emph{partially ordered set}).

L'ordine si dice \textbf{totale} (o lineare) se ogni coppia di elementi è confrontabile:
\[
    \forall x, y \in P, \quad x \le y \vee y \le x
\]
In tal caso, il poset è detto \textbf{catena}.
\end{definizione}

\subsection{Principio di Dualità per Poset}
Data una relazione d'ordine $\le$ su $P$, la relazione inversa $\ge$ (definita da $x \ge y \iff y \le x$) è anch'essa una relazione d'ordine parziale su $P$, detta \textbf{ordine duale}.
Ne discende il \textbf{Principio di Dualità}: ogni teorema valido per una relazione d'ordine $\le$ rimane valido sostituendo simultaneamente ogni occorrenza di $\le$ con $\ge$ (e scambiando i concetti duali: minimo con massimo, $\inf$ con $\sup$, minoranti con maggioranti).

\section{Elementi Speciali in un Poset}

Sia $(P, \le)$ un poset e sia $S \sseq P$.

\begin{definizione}{Minimo, Massimo, Minimali e Massimali}{elementi_speciali}
Sia $m \in P$:
\begin{itemize}
    \item $m$ è il \textbf{minimo} di $P$ ($m = \min(P)$) se:
    \[
        \forall x \in P, \quad m \le x
    \]
    \item $M$ è il \textbf{massimo} di $P$ ($M = \max(P)$) se:
    \[
        \forall x \in P, \quad x \le M
    \]
    \item $u \in P$ è un \textbf{elemento minimale} se nessun elemento di $P$ è strettamente minore di $u$:
    \[
        \forall x \in P, \quad x \le u \Implica x = u
    \]
    \item $v \in P$ è un \textbf{elemento massimale} se nessun elemento di $P$ è strettamente maggiore di $v$:
    \[
        \forall x \in P, \quad v \le x \Implica x = v
    \]
\end{itemize}
\end{definizione}

\begin{proposizione}{Unicità del Minimo e del Massimo}{unicita_min_max}
In qualunque poset $(P, \le)$, se il minimo (risp. massimo) esiste, esso è \textbf{unico}.
\end{proposizione}
\begin{proof}
Siano $m_1, m_2$ due minimi di $P$. Essendo $m_1$ un minimo e $m_2 \in P$, si ha $m_1 \le m_2$. Essendo $m_2$ un minimo e $m_1 \in P$, si ha $m_2 \le m_1$. Per l'antisimmetria della relazione d'ordine, $m_1 = m_2$.
\end{proof}

\begin{avvisoCutolo}
    \textbf{Minimo vs Elemento Minimale: Una Distinzione Cruciale!}
    \begin{itemize}
        \item Se un poset possiede un \textbf{minimo} $m$, allora $m$ è anche l'\textbf{unico elemento minimale} di $P$.
        \item Il viceversa è falso in generale per insiemi infiniti, ma è vero se il poset è finito ed esiste un solo minimale.
        \item Un poset può non avere minimo e possedere \textbf{molteplici elementi minimali} (in tal caso reciprocamente non confrontabili), oppure non averne alcuno.
    \end{itemize}
\end{avvisoCutolo}

\begin{teorema}{Esistenza di Minimali e Massimali in Poset Finiti}{teo_esist_finiti}
Ogni poset finito non vuoto $(P, \le)$ ammette \textbf{almeno un elemento minimale} e \textbf{almeno un elemento massimale}.
\end{teorema}

\subsection{Minoranti, Maggioranti, Estremo Inferiore ed Estremo Superiore}

Sia $(P, \le)$ un poset e sia $S \sseq P$:
\begin{itemize}
    \item L'insieme dei \textbf{minoranti} di $S$ è $\minor(S) = \{x \in P \mid \forall s \in S, x \le s\}$.
    \item L'insieme dei \textbf{maggioranti} di $S$ è $\maggior(S) = \{x \in P \mid \forall s \in S, s \le x\}$.
    \item Se $\minor(S)$ ammette massimo in $P$, tale massimo è detto \textbf{estremo inferiore} di $S$ e si indica con $\inf(S)$ o $\bigwedge S$:
    \[
        \inf(S) = \max(\minor(S))
    \]
    \item Se $\maggior(S)$ ammette minimo in $P$, tale minimo è detto \textbf{estremo superiore} di $S$ e si indica con $\sup(S)$ o $\bigvee S$:
    \[
        \sup(S) = \min(\maggior(S))
    \]
\end{itemize}

Nel caso di una coppia $S = \{a, b\}$, si usa la notazione $a \wedge b = \inf(a, b)$ e $a \vee b = \sup(a, b)$.

\section{Relazione di Copertura e Diagrammi di Hasse}

Nei poset finiti, gran parte della relazione d'ordine è ridondante a causa delle proprietà riflessiva e transitiva. La struttura essenziale è catturata dalla relazione di copertura.

\begin{definizione}{Relazione di Copertura}{copertura}
Dati due elementi $x, y \in P$, si dice che $y$ \textbf{copre} $x$ (e si scrive $x \prec y$) se:
\[
    x < y \quad \wedge \quad \neg(\exists z \in P : x < z < y)
\]
ovvero se non esiste alcun elemento intermedio strettamente compreso tra $x$ e $y$.
\end{definizione}

\begin{definizione}{Diagramma di Hasse}{hasse}
Il \textbf{diagramma di Hasse} di un poset finito $(P, \le)$ è una rappresentazione grafica planare in cui:
\begin{enumerate}
    \item Gli elementi di $P$ sono rappresentati da punti (o nodi) nel piano;
    \item Se $x \prec y$, il nodo $y$ è collocato a un'altezza verticale superiore rispetto a $x$;
    \item Un segmento congiunge $x$ e $y$ se e solo se $y$ copre $x$ ($x \prec y$).
\end{enumerate}
\end{definizione}

Per risalire alla relazione d'ordine a partire dal diagramma: $a \le b$ se e solo se $a = b$ oppure esiste un cammino ascendente di segmenti che sale da $a$ fino a $b$.

\begin{figure}[htbp]
\centering
\begin{tikzpicture}[scale=1.2, every node/.style={circle, draw, fill=white, inner sep=2pt}]
    % Diagramma di Hasse per P({1,2,3})
    \node (0) at (0, 0) {$\vuoto$};
    \node (1) at (-1.5, 1.2) {$\{1\}$};
    \node (2) at (0, 1.2) {$\{2\}$};
    \node (3) at (1.5, 1.2) {$\{3\}$};
    \node (12) at (-1.5, 2.5) {$\{1,2\}$};
    \node (13) at (0, 2.5) {$\{1,3\}$};
    \node (23) at (1.5, 2.5) {$\{2,3\}$};
    \node (123) at (0, 3.7) {$\{1,2,3\}$};
    
    \draw (0) -- (1); \draw (0) -- (2); \draw (0) -- (3);
    \draw (1) -- (12); \draw (1) -- (13);
    \draw (2) -- (12); \draw (2) -- (23);
    \draw (3) -- (13); \draw (3) -- (23);
    \draw (12) -- (123); \draw (13) -- (123); \draw (23) -- (123);
\end{tikzpicture}
\caption{Diagramma di Hasse del reticolo booleano $(\P(\{1,2,3\}), \sseq)$.}
\label{fig:hasse_p3}
\end{figure}

\section{Reticoli: Definizione d'Ordine ed Algebrica}

\begin{definizione}{Reticolo (come Poset)}{reticolo_poset}
Un poset $(L, \le)$ si dice \textbf{reticolo} se per ogni coppia di elementi $a, b \in L$ esistono in $L$ sia l'estremo inferiore che l'estremo superiore:
\[
    \forall a, b \in L, \quad \exists (a \wedge b) = \inf(a, b) \quad \wedge \quad \exists (a \vee b) = \sup(a, b)
\]
\end{definizione}

\begin{definizione}{Reticolo (come Struttura Algebrica)}{reticolo_alg}
Una struttura algebrica $(L, \wedge, \vee)$ dotata di due operazioni binarie interne si dice \textbf{reticolo} se per ogni $a, b, c \in L$ sono soddisfatti i seguenti assiomi:
\begin{enumerate}
    \item \textbf{Commutatività:} $a \wedge b = b \wedge a$, \quad $a \vee b = b \vee a$;
    \item \textbf{Associatività:} $(a \wedge b) \wedge c = a \wedge (b \wedge c)$, \quad $(a \vee b) \vee c = a \vee (b \vee c)$;
    \item \textbf{Idempotenza:} $a \wedge a = a$, \quad $a \vee a = a$;
    \item \textbf{Leggi di Assorbimento:} $a \wedge (a \vee b) = a$, \quad $a \vee (a \wedge b) = a$.
\end{enumerate}
\end{definizione}

\begin{teorema}{Teorema di Equivalenza tra le Definizioni di Reticolo}{teo_equiv_reticolo}
Le nozioni di reticolo come poset e di reticolo come struttura algebrica sono \textbf{perfettamente equivalenti}:
\begin{enumerate}
    \item Se $(L, \le)$ è un reticolo d'ordine, ponendo $a \wedge b = \inf(a, b)$ e $a \vee b = \sup(a, b)$, la struttura $(L, \wedge, \vee)$ soddisfa i 4 assiomi algebrici.
    \item Viceversa, se $(L, \wedge, \vee)$ è un reticolo algebrico, la relazione binaria definita su $L$ da:
    \[
        a \le b \quad \Shiff \quad a \wedge b = a \quad \Shiff \quad a \vee b = b
    \]
    è una relazione d'ordine parziale su $L$ tale che $\inf(a, b) = a \wedge b$ e $\sup(a, b) = a \vee b$.
\end{enumerate}
\end{teorema}

\begin{avvisoCutolo}
    \textbf{Dimostrazione Obbligatoria all'Orale: L'Associatività di $\inf$!}
    Il Prof. Cutolo chiede sistematicamente la dimostrazione formale che $\inf$ è associativo \textbf{lavorando unicamente con la definizione d'ordine di estremo inferiore}, senza assumere a priori l'assioma algebrico!
\end{avvisoCutolo}

\begin{proof}[Dimostrazione che $\inf$ è associativo in un poset]
Sia $(L, \le)$ un reticolo d'ordine e siano $a, b, c \in L$.
Poniamo $d_1 = a \wedge (b \wedge c)$ e $d_2 = (a \wedge b) \wedge c$. Vogliamo dimostrare che $d_1 = d_2$.
Mostriamo anzitutto che $d_1 = \inf(\{a, b, c\})$.
Per definizione di estremo inferiore:
\begin{enumerate}
    \item $d_1 = \inf(a, b \wedge c) \implies d_1 \le a$ e $d_1 \le (b \wedge c)$.
    \item A sua volta, $d_1 \le (b \wedge c) \implies d_1 \le b$ e $d_1 \le c$.
    Dunque $d_1$ è un minorante comune di $\{a, b, c\}$, cioè $d_1 \in \minor(\{a, b, c\})$.
    \item Sia ora $x \in L$ un qualunque minorante comune di $\{a, b, c\}$, ovvero $x \le a$, $x \le b$ e $x \le c$.
    Poiché $x \le b$ e $x \le c$, per la proprietà di massimo dei minoranti si ha $x \le \inf(b, c) = b \wedge c$.
    Ora abbiamo $x \le a$ e $x \le (b \wedge c)$. Di conseguenza, $x \le \inf(a, b \wedge c) = d_1$.
\end{enumerate}
Ciò prova che $d_1$ è il massimo dei minoranti di $\{a, b, c\}$, cioè $d_1 = \inf(\{a, b, c\})$.
Ripetendo lo stesso identico ragionamento per $d_2 = (a \wedge b) \wedge c$, si ottiene che anch'esso è l'estremo inferiore di $\{a, b, c\}$.
Poiché per la Proposizione~\ref{prop:unicita_min_max} l'estremo inferiore (essendo un massimo) è unico se esiste, si conclude che $d_1 = d_2$.
\end{proof}

\subsection{Reticoli Limitati e Reticoli Completi}
Un reticolo $(L, \wedge, \vee)$ si dice:
\begin{itemize}
    \item \textbf{Limitato} se possiede sia elemento minimo $0 = \min(L)$ che massimo $1 = \max(L)$. In tal caso:
    \[
        \forall x \in L, \quad x \wedge 0 = 0, \quad x \vee 0 = x, \quad x \wedge 1 = x, \quad x \vee 1 = 1
    \]
    Ogni reticolo finito non vuoto è limitato: $0 = \bigwedge L$ e $1 = \bigvee L$.
    \item \textbf{Completo} se ogni sottoinsieme non vuoto $S \sseq L$ (anche infinito) ammette estremo inferiore e superiore in $L$.
\end{itemize}

\section{Reticoli Distributivi e Criterio di Birkhoff}

\begin{definizione}{Reticolo Distributivo}{ret_distributivo}
Un reticolo $(L, \wedge, \vee)$ si dice \textbf{distributivo} se l'operazione $\wedge$ si distribuisce su $\vee$:
\[
    \forall a, b, c \in L, \quad a \wedge (b \vee c) = (a \wedge b) \vee (a \wedge c)
\]
In ogni reticolo, la validità di questa identità è \textbf{equivalente} alla validità della legge distributiva duale:
\[
    a \vee (b \wedge c) = (a \vee b) \wedge (a \vee c)
\]
\end{definizione}

In ogni reticolo vale sempre la disuguaglianza distributiva:
\[
    (a \wedge b) \vee (a \wedge c) \le a \wedge (b \vee c)
\]
Nei reticoli non distributivi tale disuguaglianza può essere stretta.

\begin{teorema}{Criterio di Distributività di Birkhoff}{birkhoff}
Un reticolo $L$ è distributivo \textbf{se e solo se} non contiene alcun sottoreticolo isomorfo al \textbf{pentagono} $N_5$ o al \textbf{diamante} $M_3$.
\end{teorema}

\begin{figure}[htbp]
\centering
\begin{minipage}{0.45\textwidth}
\centering
\begin{tikzpicture}[scale=1.1, every node/.style={circle, draw, fill=white, inner sep=2pt}]
    \node (0) at (0, 0) {$0$};
    \node (a) at (-1, 1.5) {$a$};
    \node (b) at (1, 1) {$b$};
    \node (c) at (1, 2) {$c$};
    \node (1) at (0, 3) {$1$};
    \draw (0) -- (a) -- (1);
    \draw (0) -- (b) -- (c) -- (1);
\end{tikzpicture}
\caption*{Pentagono $N_5$ (Non modulare)}
\end{minipage}
\hfill
\begin{minipage}{0.45\textwidth}
\centering
\begin{tikzpicture}[scale=1.1, every node/.style={circle, draw, fill=white, inner sep=2pt}]
    \node (0) at (0, 0) {$0$};
    \node (a) at (-1.2, 1.5) {$a$};
    \node (b) at (0, 1.5) {$b$};
    \node (c) at (1.2, 1.5) {$c$};
    \node (1) at (0, 3) {$1$};
    \draw (0) -- (a) -- (1);
    \draw (0) -- (b) -- (1);
    \draw (0) -- (c) -- (1);
\end{tikzpicture}
\caption*{Diamante $M_3$ (Modulare non distributivo)}
\end{minipage}
\caption{I due reticoli proibiti del Criterio di Birkhoff.}
\label{fig:n5_m3}
\end{figure}

\begin{avvisoCutolo}
    Nel pentagono $N_5$, si osservi che $b < c$, ma:
    \[
        a \wedge (b \vee c) = a \wedge 1 = a \neq b = b \vee 0 = (a \wedge b) \vee (a \wedge c)
    \]
    Nel diamante $M_3$:
    \[
        a \wedge (b \vee c) = a \wedge 1 = a \neq 0 = 0 \vee 0 = (a \wedge b) \vee (a \wedge c)
    \]
    All'orale o nello scritto, per dimostrare che un reticolo non è distributivo, basta esibire una terna di elementi che viola la legge distributiva!
\end{avvisoCutolo}

\section{Complementi, Algebre di Boole e Anelli Booleani}

\begin{definizione}{Complemento e Reticolo Complementato}{complemento}
Sia $(L, \wedge, \vee, 0, 1)$ un reticolo limitato. Sia $a \in L$. Un elemento $b \in L$ si dice \textbf{complemento} di $a$ se:
\[
    a \wedge b = 0 \quad \wedge \quad a \vee b = 1
\]
Un reticolo limitato in cui \textbf{ogni elemento ammette almeno un complemento} si dice \textbf{reticolo complementato}.
\end{definizione}

\begin{teorema}{Unicità del Complemento nei Reticoli Distributivi}{unicita_comp}
Sia $L$ un reticolo \textbf{distributivo} limitato. Se un elemento $a \in L$ ammette un complemento, questo è \textbf{unico}.
\end{teorema}
\begin{proof}
Siano $b_1, b_2$ due complementi di $a$, dunque $a \wedge b_1 = 0$, $a \vee b_1 = 1$, $a \wedge b_2 = 0$, $a \vee b_2 = 1$.
Calcoliamo:
\[
    b_1 = b_1 \wedge 1 = b_1 \wedge (a \vee b_2) = (b_1 \wedge a) \vee (b_1 \wedge b_2) = 0 \vee (b_1 \wedge b_2) = b_1 \wedge b_2
\]
Dunque $b_1 \le b_2$. Scambiando i ruoli di $b_1$ e $b_2$, si ottiene analogamente $b_2 \le b_1$. Per antisimmetria, $b_1 = b_2$.
\end{proof}

\begin{definizione}{Reticolo Booleano e Algebra di Boole}{boole}
Un reticolo si dice \textbf{booleano} se è distributivo e complementato.
L'unico complemento di $a$ si denota con $a'$ (oppure $\neg a$ o $\bar{a}$).
La struttura $(B, \wedge, \vee, ', 0, 1)$ è detta \textbf{algebra di Boole}.
\end{definizione}

Nelle algebre di Boole valgono le leggi di De Morgan per il complemento:
\[
    (a \wedge b)' = a' \vee b', \qquad (a \vee b)' = a' \wedge b'
\]

\begin{teorema}{Teorema di Stone per Algebre di Boole Finite}{stone}
Ogni algebra di Boole finita $B$ è isomorfa al reticolo delle parti $(\P(S), \cap, \cup)$ dell'insieme dei suoi atomi $S = \operatorname{Atomi}(B)$.
In particolare, la cardinalità di ogni algebra di Boole finita è una \textbf{potenza di 2}:
\[
    |B| = 2^n \quad (\text{con } n = |S|)
\]
\end{teorema}

\begin{definizione}{Anello Booleano}{anello_boole}
Un anello $(R, +, \cdot)$ con unità si dice \textbf{anello booleano} se ogni suo elemento è \textbf{idempotente}:
\[
    \forall x \in R, \quad x^2 = x
\]
\end{definizione}

\begin{teorema}{Proprietà degli Anelli Booleani}{prop_anelli_boole}
Sia $R$ un anello booleano. Allora:
\begin{enumerate}
    \item $R$ ha caratteristica 2: per ogni $x \in R$, $x + x = 0$ (ovvero $x = -x$);
    \item $R$ è commutativo: per ogni $x, y \in R$, $xy = yx$.
\end{enumerate}
\end{teorema}
\begin{proof}
1. Consideriamo $(x + x)^2 = x + x$. Per distributività, $(x+x)^2 = x^2 + x^2 + x^2 + x^2 = x + x + x + x$. Uguagliando: $x + x + x + x = x + x \implies x + x = 0$.
2. Consideriamo $(x + y)^2 = x + y$. Espandendo: $x^2 + xy + yx + y^2 = x + xy + yx + y$. Poiché $x^2 = x$ e $y^2 = y$, otteniamo $xy + yx = 0$. Essendo $yx = -yx$, si ha $xy = -yx = yx$.
\end{proof}

\begin{teorema}{Equivalenza tra Algebre di Boole e Anelli Booleani}{equiv_boole}
Ogni algebra di Boole $(B, \wedge, \vee, ', 0, 1)$ diviene un anello booleano ponendo:
\[
    a + b \coloneqq (a \wedge b') \vee (a' \wedge b) = a \ds b, \qquad a \cdot b \coloneqq a \wedge b
\]
Viceversa, ogni anello booleano unitario $(R, +, \cdot)$ diviene un'algebra di Boole ponendo:
\[
    a \wedge b \coloneqq a \cdot b, \qquad a \vee b \coloneqq a + b + ab, \qquad a' \coloneqq 1 + a
\]
\end{teorema}


% ------------------------------------------------------------------------------
% PARTE III: TEORIA DEI NUMERI ED ARITMETICA MODULARE
% ------------------------------------------------------------------------------
\part{Aritmetica, Congruenze e Teoria dei Numeri}
\chapter{Aritmetica Modulare, Teorema di Bézout ed Equazioni Diofantee}
\label{chap:aritmetica}

L'aritmetica nell'anello $\Z$ e nei suoi quozienti modulari $\Z_m$ è al centro di quasi ogni prova scritta del Prof. Cutolo. In questo capitolo presentiamo l'algoritmo euclideo esteso, la risoluzione formale delle equazioni congruenziali, i sistemi cinesi del resto e i teoremi di Eulero e Fermat.

\section{Divisione con Resto e Divisibilità in $\Z$}

\begin{teorema}{Teorema della Divisione Aritmetica con Resto (Divisione Euclidea)}{div_resto}
Siano $a, b \in \Z$ con $b \neq 0$. Esiste un'\textbf{unica} coppia di interi $q, r \in \Z$ (detti rispettivamente \textbf{quoziente} e \textbf{resto}) tali che:
\[
    a = bq + r \qquad \text{con } 0 \le r < |b|
\]
\end{teorema}

\begin{proof}
Dimostriamo esistenza e unicità:
\begin{enumerate}
    \item \emph{Esistenza:} Consideriamo l'insieme $S = \{a - bx \mid x \in \Z, a - bx \ge 0\}$.
    $S$ è un sottoinsieme non vuoto di $\N$. Per il Principio del Buon Ordinamento (Teorema~\ref{teo:buon_ord}), $S$ ammette un minimo elemento $r = a - bq \ge 0$ per un opportuno $q \in \Z$.
    Mostriamo che $r < |b|$. Se per assurdo fosse $r \ge |b|$, potremmo considerare $r' = r - |b| = a - b(q \pm 1) \ge 0$. Allora $r' \in S$ con $r' < r$, contraddicendo la minimalità di $r$ in $S$. Dunque $0 \le r < |b|$.
    \item \emph{Unicità:} Siano $(q, r)$ e $(q', r')$ due coppie soddisfacenti le condizioni: $a = bq + r = bq' + r'$ con $0 \le r, r' < |b|$.
    Allora $b(q - q') = r' - r$. Prendendo i valori assoluti: $|b| \cdot |q - q'| = |r' - r|$.
    Poiché $0 \le r, r' < |b|$, la loro distanza deve soddisfare $|r' - r| < |b|$.
    Ne segue che $|b| \cdot |q - q'| < |b|$, il che è possibile solo se $|q - q'| = 0$, da cui $q = q'$ e di conseguenza $r = r'$.
\end{enumerate}
\end{proof}

\begin{definizione}{Divisibilità in $\Z$}{divisibilita}
Dati due interi $a, b \in \Z$, diciamo che $a$ \textbf{divide} $b$ (e scriviamo $a \mid b$) se esiste un intero $c \in \Z$ tale che $b = ac$.
Gli elementi invertibili di $\Z$ sono $\mathcal{U}(\Z) = \{1, -1\}$.
Due interi $a, b$ si dicono \textbf{associati} in $\Z$ ($a \sim b$) se $a \mid b$ e $b \mid a$, ovvero se $a = \pm b$.
\end{definizione}

\section{Massimo Comune Divisore e Teorema di Bézout}

\begin{definizione}{Massimo Comune Divisore}{mcd}
Siano $a, b \in \Z$ non entrambi nulli. Un intero $d \ge 1$ è detto \textbf{massimo comune divisore} di $a$ e $b$ (e si denota con $\mcd(a, b)$ o $\gcd(a, b)$) se:
\begin{enumerate}
    \item $d$ è un divisore comune: $d \mid a$ e $d \mid b$;
    \item ogni divisore comune di $a$ e $b$ divide $d$: per ogni $c \in \Z$, se $c \mid a$ e $c \mid b$, allora $c \mid d$.
\end{enumerate}
Se $\mcd(a, b) = 1$, $a$ e $b$ si dicono \textbf{coprimi} (o primi tra loro).
\end{definizione}

\begin{teorema}{Teorema di Bézout}{bezout}
Siano $a, b \in \Z$ non entrambi nulli e sia $d = \mcd(a, b)$. Allora esistono due interi $u, v \in \Z$ (detti \textbf{coefficienti di Bézout}) tali che:
\[
    d = au + bv
\]
Inoltre, l'insieme delle combinazioni lineari a coefficienti interi di $a$ e $b$ coincide esattamente con l'insieme di tutti i multipli del loro massimo comune divisore:
\[
    \{ax + by \mid x, y \in \Z\} = d\Z
\]
\end{teorema}

\subsection{L'Algoritmo Euclideo Esteso (Metodo a Tabella di Cutolo)}
L'algoritmo euclideo calcola $\mcd(a, b)$ eseguendo divisioni con resto successive fino ad ottenere resto nullo. L'ultimo resto non nullo è esattamente $\mcd(a, b)$.
Sia $a > b > 0$:
\begin{align*}
    a &= b q_1 + r_1 && 0 < r_1 < b \\
    b &= r_1 q_2 + r_2 && 0 < r_2 < r_1 \\
    r_1 &= r_2 q_3 + r_3 && 0 < r_3 < r_2 \\
    &\;\;\vdots && \\
    r_{k-2} &= r_{k-1} q_k + r_k && r_k = \mcd(a, b) \\
    r_{k-1} &= r_k q_{k+1} + 0 && (\text{resto nullo})
\end{align*}

Per ricavare i coefficienti $u, v$ di Bézout si procede a ritroso oppure si compila la \textbf{tabella matriciale} passo per passo:
\begin{center}
\begin{tabular}{|c|c|c|c|}
\hline
$i$ & Resto $r_i$ & $u_i$ (coord. $a$) & $v_i$ (coord. $b$) \\
\hline
$-1$ & $a$ & $1$ & $0$ \\
$0$ & $b$ & $0$ & $1$ \\
$1$ & $r_1 = a - q_1 b$ & $1$ & $-q_1$ \\
$\vdots$ & $r_i = r_{i-2} - q_i r_{i-1}$ & $u_i = u_{i-2} - q_i u_{i-1}$ & $v_i = v_{i-2} - q_i v_{i-1}$ \\
\hline
\end{tabular}
\end{center}

\begin{teorema}{Lemma di Euclide e Teorema Fondamentale dell'Aritmetica}{fond_aritmetica}
Un intero $p > 1$ è detto \textbf{primo} se i suoi unici divisori positivi sono $1$ e $p$.
\begin{enumerate}
    \item \textbf{Lemma di Euclide:} se $p$ è primo e $p \mid ab$, allora $p \mid a$ oppure $p \mid b$.
    Più in generale: se $\mcd(a, b) = 1$ e $a \mid bc$, allora $a \mid c$.
    \item \textbf{Teorema Fondamentale dell'Aritmetica:} ogni intero $n \ge 2$ si fattorizza in modo unico (a meno dell'ordine dei fattori) come prodotto di potenze di numeri primi distinti:
    \[
        n = p_1^{e_1} p_2^{e_2} \cdots p_k^{e_k} \qquad (p_i \text{ primi distinti}, e_i \ge 1)
    \]
\end{enumerate}
\end{teorema}

\section{Congruenze Modulari e l'Anello $\Z_m$}

\begin{definizione}{Relazione di Congruenza Modulo $m$}{congruenza}
Sia $m \in \N^*$ fissato. Due interi $a, b \in \Z$ si dicono \textbf{congrui modulo $m$} (scritto $a \equiv b \pmod m$) se la loro differenza è multiplo di $m$:
\[
    a \equiv b \pmod m \quad \Shiff \quad m \mid (a - b) \quad \Shiff \quad a \text{ e } b \text{ hanno lo stesso resto diviso } m
\]
\end{definizione}

La congruenza modulo $m$ è una \textbf{relazione di equivalenza} su $\Z$. Inoltre, è una \textbf{congruenza algebrica}, ovvero compatibile con la somma e con il prodotto in $\Z$:
\[
    a_1 \equiv b_1 \pmod m \;\wedge\; a_2 \equiv b_2 \pmod m \Implica \begin{cases}
        a_1 + a_2 \equiv b_1 + b_2 \pmod m \\
        a_1 \cdot a_2 \equiv b_1 \cdot b_2 \pmod m
    \end{cases}
\]

\begin{definizione}{L'Anello Quoziente $\Z_m$}{anello_zm}
L'insieme quoziente $\Z/\equiv_m$ è denotato con $\Z_m$ (oppure $\Z/m\Z$) ed è costituito da esattamente $m$ classi di equivalenza distinte:
\[
    \Z_m = \{[0]_m, [1]_m, [2]_m, \dots, [m-1]_m\}
\]
Dotato delle operazioni quoziente $[a] + [b] = [a + b]$ e $[a] \cdot [b] = [ab]$, $\Z_m$ è un \textbf{anello commutativo unitario}.
\end{definizione}

\begin{teorema}{Elementi Invertibili e Divisori dello Zero in $\Z_m$}{inv_zm}
Sia $m \ge 2$ e sia $[a] \in \Z_m$.
\begin{enumerate}
    \item $[a]$ è un \textbf{elemento invertibile} di $\Z_m$ se e solo se $\mcd(a, m) = 1$:
    \[
        \mathcal{U}(\Z_m) = \{[a] \in \Z_m \mid \mcd(a, m) = 1\}
    \]
    L'inverso aritmetico $[a]^{-1} = [u]$ è dato dal coefficiente di Bézout $au + mv = 1$.
    \item $[a]$ è \textbf{cancellabile} in $\Z_m$ se e solo se è invertibile.
    \item $[a]$ è un \textbf{divisore dello zero} di $\Z_m$ se e solo se $[a] \neq [0]$ e $\mcd(a, m) > 1$.
    \item L'anello $\Z_m$ è un \textbf{campo} (e un dominio d'integrità) se e solo se $m$ è un numero \textbf{primo}.
\end{enumerate}
\end{teorema}

\section{Risoluzione di Equazioni Congruenziali Lineari}

Un'equazione congruenziale lineare a una incognita ha la forma:
\begin{equation}
\label{eq:congruenza_lin}
    ax \equiv b \pmod m
\end{equation}
Risolvere l'equazione~\eqref{eq:congruenza_lin} equivale a determinare le soluzioni intere $(x, y)$ dell'\textbf{equazione diofantea lineare}:
\[
    ax - my = b
\]

\begin{teorema}{Criterio di Solvibilità per Equazioni Congruenziali Lineari}{solv_congr}
Sia $d = \mcd(a, m)$.
\begin{enumerate}
    \item L'equazione $ax \equiv b \pmod m$ ammette soluzioni se e solo se $d$ divide $b$:
    \[
        d \mid b
    \]
    Se $d \nmid b$, l'equazione è \textbf{impossibile} (non ha soluzioni).
    \item Se $d \mid b$, l'equazione ammette infinite soluzioni in $\Z$, le quali formano esattamente \textbf{$d$ classi di equivalenza distinte modulo $m$}.
\end{enumerate}
\end{teorema}

\subsection{Procedura di Risoluzione Passo per Passo}
Data l'equazione $ax \equiv b \pmod m$:
\begin{enumerate}
    \item \textbf{Calcolo del MCD:} Si calcola $d = \mcd(a, m)$ con l'algoritmo euclideo. Se $d \nmid b$, stop (nessuna soluzione).
    \item \textbf{Semplificazione:} Se $d \mid b$, si divide l'intera congruenza per $d$:
    \[
        a' x \equiv b' \pmod{m'}, \qquad \text{ove } a' = \frac{a}{d}, \quad b' = \frac{b}{d}, \quad m' = \frac{m}{d}
    \]
    Per costruzione, $\mcd(a', m') = 1$, dunque $[a']$ è invertibile modulo $m'$.
    \item \textbf{Ricerca dell'inverso:} Tramite l'algoritmo di Bézout esteso si trovano $u, v \in \Z$ tali che $a'u + m'v = 1$, da cui $a'u \equiv 1 \pmod{m'}$. L'inverso di $a'$ modulo $m'$ è $u$.
    \item \textbf{Soluzione particolare:} Moltiplicando per $u$:
    \[
        x_0 \equiv u \cdot b' \pmod{m'}
    \]
    \item \textbf{Tutte le $d$ soluzioni distinte modulo $m$:}
    \[
        x_k = x_0 + k \cdot m' = x_0 + k \cdot \frac{m}{d}, \qquad \text{per } k = 0, 1, 2, \dots, d-1
    \]
\end{enumerate}

\section{Sistemi di Congruenze e Teorema Cinese del Resto (CRT)}

Consideriamo un sistema di $k$ equazioni congruenziali lineari:
\[
\begin{cases}
    x \equiv c_1 \pmod{m_1} \\
    x \equiv c_2 \pmod{m_2} \\
    \quad \vdots \\
    x \equiv c_k \pmod{m_k}
\end{cases}
\]

\begin{teorema}{Teorema Cinese del Resto (Chinese Remainder Theorem)}{crt}
Siano $m_1, m_2, \dots, m_k \in \N_{\ge 2}$ moduli \textbf{a due a due coprimi}, ovvero:
\[
    \mcd(m_i, m_j) = 1 \quad \text{per ogni } i \neq j
\]
Allora, per qualunque scelta dei resti $c_1, c_2, \dots, c_k \in \Z$, il sistema ammette soluzione in $\Z$.
Tale soluzione è \textbf{unica modulo il prodotto dei moduli} $M = m_1 m_2 \cdots m_k$.
\end{teorema}

\begin{proof}[Dimostrazione Costruttiva]
Poniamo $M = m_1 m_2 \cdots m_k$ e per ciascun $i \in \{1, \dots, k\}$ definiamo:
\[
    M_i = \frac{M}{m_i} = \prod_{j \neq i} m_j
\]
Poiché i moduli sono a due a due coprimi, si ha $\mcd(M_i, m_i) = 1$.
Dunque per il Teorema~\ref{teo:inv_zm}, ciascun $M_i$ è invertibile modulo $m_i$. Esiste quindi un intero $N_i \in \Z$ tale che:
\[
    M_i N_i \equiv 1 \pmod{m_i}
\]
Definiamo la combinazione:
\[
    x_0 = \sum_{i=1}^k c_i M_i N_i
\]
Verifichiamo che $x_0$ soddisfa la $j$-esima congruenza del sistema:
Per $i \neq j$, $M_i$ contiene il fattore $m_j$, quindi $M_i \equiv 0 \pmod{m_j}$.
Dunque nella sommatoria tutti i termini con $i \neq j$ si annullano modulo $m_j$:
\[
    x_0 \equiv c_j M_j N_j \equiv c_j \cdot 1 = c_j \pmod{m_j}
\]
Ciò vale per ogni $j \in \{1, \dots, k\}$. Dunque $x_0$ è soluzione del sistema. Tutte le soluzioni sono date da $x \equiv x_0 \pmod M$.
\end{proof}

\subsection{Sistemi con Moduli non Coprimi}
Se $\mcd(m_i, m_j) = d > 1$, il sistema ammette soluzione se e solo se per ogni coppia vale la condizione di compatibilità $c_i \equiv c_j \pmod{\mcd(m_i, m_j)}$.
In tal caso si decompone ciascuna congruenza nelle sue componenti prime mediante fattorizzazione.

\section{La Funzione $\varphi$ di Eulero e i Teoremi di Eulero e Fermat}

\begin{definizione}{Funzione $\varphi$ di Eulero}{eulerphi}
La \textbf{funzione $\varphi$ di Eulero} (o indicatore di Eulero) $\varphi: \N^* \to \N^*$ conta quanti numeri interi positivi minori o uguali a $m$ sono coprimi con $m$:
\[
    \varphi(m) = |\{a \in \{1, 2, \dots, m\} \mid \mcd(a, m) = 1\}| = |\mathcal{U}(\Z_m)|
\]
\end{definizione}

\begin{teorema}{Calcolo della Funzione $\varphi$}{calcolo_phi}
La funzione $\varphi$ gode delle seguenti proprietà:
\begin{enumerate}
    \item Se $p$ è un numero primo, $\varphi(p) = p - 1$.
    \item Per ogni potenza di un primo $p^k$ ($k \ge 1$):
    \[
        \varphi(p^k) = p^k - p^{k-1} = p^k \left(1 - \frac{1}{p}\right)
    \]
    \item È \textbf{moltiplicativa}: se $\mcd(m, n) = 1$, allora $\varphi(mn) = \varphi(m)\varphi(n)$.
    \item In generale, se $m = p_1^{e_1} p_2^{e_2} \cdots p_k^{e_k}$, allora:
    \[
        \varphi(m) = m \prod_{i=1}^k \left(1 - \frac{1}{p_i}\right) = \prod_{i=1}^k p_i^{e_i - 1}(p_i - 1)
    \]
\end{enumerate}
\end{teorema}

\begin{teorema}{Teorema di Eulero-Fermat e Piccolo Teorema di Fermat}{teoremi_fermat}
\begin{enumerate}
    \item \textbf{Teorema di Eulero-Fermat:} Per ogni intero $a \in \Z$ coprimo con $m$ ($\mcd(a, m) = 1$):
    \[
        a^{\varphi(m)} \equiv 1 \pmod m
    \]
    \item \textbf{Piccolo Teorema di Fermat:} Se $p$ è un numero primo e $p \nmid a$, allora:
    \[
        a^{p-1} \equiv 1 \pmod p
    \]
    Inoltre, per qualunque intero $a \in \Z$ vale $a^p \equiv a \pmod p$.
\end{enumerate}
\end{teorema}

\subsection{Applicazione al Calcolo di Grandi Potenze Modulari}
Per calcolare il resto della divisione di $a^N$ per $m$ con $\mcd(a, m) = 1$:
\begin{enumerate}
    \item Si calcola $\varphi(m)$.
    \item Si divide l'esponente $N$ per $\varphi(m)$: $N = q \cdot \varphi(m) + r$ con $0 \le r < \varphi(m)$.
    \item Si applica il teorema: $a^N = (a^{\varphi(m)})^q \cdot a^r \equiv 1^q \cdot a^r = a^r \pmod m$.
\end{enumerate}
Si riduce così un esponente astronomico a un esponente $r < \varphi(m)$ facilmente calcolabile!


% ------------------------------------------------------------------------------
% PARTE IV: STRUTTURE ALGEBRICHE E POLINOMI
% ------------------------------------------------------------------------------
\part{Strutture Algebriche, Gruppi, Anelli e Polinomi}
\chapter{Operazioni Interne e Strutture Algebriche Fondamentali}
\label{chap:strutture_algebriche}

Le strutture algebriche costituiscono l'oggetto centrale dell'Algebra Astratta. In questo capitolo studiamo assiomaticamente le operazioni binarie interne, gli elementi notevoli (neutri, simmetrici, cancellabili), i semigruppi, i monoidi, i gruppi, gli anelli, i domini d'integrità e i campi.

\section{Operazioni Binarie Interne ed Assiomi}

\begin{definizione}{Operazione Binaria Interna}{operazione}
Sia $S$ un insieme. Un'\textbf{operazione binaria interna} su $S$ è un'applicazione:
\[
    * : S \times S \longrightarrow S, \qquad (a, b) \longmapsto a * b
\]
Un sottoinsieme non vuoto $T \sseq S$ si dice \textbf{chiuso} (o \textbf{stabile}) rispetto a $*$ se:
\[
    \forall a, b \in T, \quad a * b \in T
\]
In tal caso, la restrizione $*|_T : T \times T \to T$ è detta \textbf{operazione indotta} su $T$.
\end{definizione}

\subsection{Proprietà delle Operazioni}
L'operazione $*$ su $S$ può godere delle seguenti proprietà:
\begin{enumerate}
    \item \textbf{Associatività:} $\forall a, b, c \in S, \quad (a * b) * c = a * (b * c)$.
    In virtù del \emph{Teorema di Associatività Generalizzata}, quando un'operazione è associativa è possibile omettere le parentesi per qualunque numero di operandi: $a_1 * a_2 * \dots * a_n$.
    \item \textbf{Commutatività:} $\forall a, b \in S, \quad a * b = b * a$.
    \item \textbf{Elemento Neutro:}
    \begin{itemize}
        \item $e_d \in S$ è un \textbf{neutro a destra} se $\forall a \in S, a * e_d = a$.
        \item $e_s \in S$ è un \textbf{neutro a sinistra} se $\forall a \in S, e_s * a = a$.
        \item $e \in S$ è un \textbf{neutro (bilatero)} se è simultaneamente neutro a destra e a sinistra:
        \[
            \forall a \in S, \quad a * e = e * a = a
        \]
    \end{itemize}
\end{enumerate}

\begin{teorema}{Unicità dell'Elemento Neutro}{unicita_neutro}
Se in una struttura $(S, *)$ esistono un elemento neutro a sinistra $e_s$ e un elemento neutro a destra $e_d$, allora essi \textbf{coincidono}:
\[
    e_s = e_d = e
\]
In particolare, se una struttura possiede un elemento neutro bilatero, questo è \textbf{unico}.
\end{teorema}
\begin{proof}
Calcoliamo l'espressione $e_s * e_d$:
\begin{itemize}
    \item Poiché $e_s$ è neutro a sinistra: $e_s * e_d = e_d$.
    \item Poiché $e_d$ è neutro a destra: $e_s * e_d = e_s$.
\end{itemize}
Uguagliando i due risultati: $e_s = e_d$.
\end{proof}

\begin{avvisoCutolo}
    Possono esistere più neutri a destra in assenza di un neutro a sinistra!
    \emph{Controesempio classico di Cutolo:} Si consideri un insieme $S$ con almeno due elementi con l'operazione proiettiva $a * b \coloneqq a$.
    Allora per ogni elemento $x \in S$ si ha $a * x = a$, dunque \textbf{ogni elemento di $S$ è un neutro a destra}! Tuttavia non esiste alcun neutro a sinistra, poiché $x * a = x \neq a$ se $x \neq a$.
\end{avvisoCutolo}

\subsection{Elementi Simmetrici (Inversi)}
Sia $(S, *)$ una struttura dotata di elemento neutro bilatero $e$. Sia $a \in S$:
\begin{itemize}
    \item $a'_d \in S$ è un \textbf{simmetrico destro} di $a$ se $a * a'_d = e$.
    \item $a'_s \in S$ è un \textbf{simmetrico sinistro} di $a$ se $a'_s * a = e$.
    \item $a' \in S$ è un \textbf{simmetrico (o inverso bilatero)} di $a$ se $a * a' = a' * a = e$.
    Un elemento che ammette simmetrico bilatero si dice \textbf{simmetrizzabile} (o \textbf{invertibile}).
\end{itemize}

\begin{teorema}{Unicità del Simmetrico sotto Associatività}{unicita_simmetrico}
Sia $(S, *)$ una struttura con operazione \textbf{associativa} e dotata di elemento neutro $e$.
Se un elemento $a \in S$ ammette un simmetrico sinistro $a'_s$ e un simmetrico destro $a'_d$, allora essi \textbf{coincidono}:
\[
    a'_s = a'_d = a'
\]
Di conseguenza, il simmetrico bilatero (se esiste) è \textbf{unico} e viene denotato con $a^{-1}$ (nella notazione moltiplicativa) o con $-a$ (nella notazione additiva).
\end{teorema}
\begin{proof}
Calcoliamo il prodotto triplo $(a'_s * a) * a'_d$:
\[
    a'_s = a'_s * e = a'_s * (a * a'_d) \overset{\text{assoc.}}{=} (a'_s * a) * a'_d = e * a'_d = a'_d
\]
\end{proof}

\subsection{Elementi Cancellabili}

\begin{definizione}{Elemento Cancellabile}{cancellabile}
Sia $(S, *)$ una struttura algebrica. Un elemento $a \in S$ si dice:
\begin{itemize}
    \item \textbf{Cancellabile a sinistra} se l'applicazione traslazione sinistra $\lambda_a: S \to S$, $\lambda_a(x) = a * x$ è iniettiva:
    \[
        \forall x, y \in S, \quad a * x = a * y \Implica x = y
    \]
    \item \textbf{Cancellabile a destra} se la traslazione destra $\rho_a: S \to S$, $\rho_a(x) = x * a$ è iniettiva:
    \[
        \forall x, y \in S, \quad x * a = y * a \Implica x = y
    \]
    \item \textbf{Cancellabile (bilatero)} se è cancellabile sia a destra che a sinistra.
\end{itemize}
\end{definizione}

\begin{teorema}{Simmetrizzabilità e Cancellabilità}{simm_canc}
Sia $(M, *, e)$ un monoide (operazione associativa con neutro).
\begin{enumerate}
    \item Ogni elemento \textbf{simmetrizzabile} è \textbf{cancellabile}:
    \[
        a \in \mathcal{U}(M) \Implica a \text{ è cancellabile}
    \]
    \item Il viceversa è falso in generale per monoidi infiniti (ad es. in $(\Z, \cdot)$, $2$ è cancellabile ma non invertibile).
    \item Se $M$ è un \textbf{monoide finito}, allora ogni elemento cancellabile è \textbf{invertibile}:
    \[
        M \text{ finito } \Implica (a \text{ è cancellabile } \iff a \text{ è invertibile})
    \]
\end{enumerate}
\end{teorema}
\begin{proof}
1. Se $a$ ha inverso $a^{-1}$, allora $a * x = a * y \implies a^{-1} * (a * x) = a^{-1} * (a * y) \implies (a^{-1} * a) * x = (a^{-1} * a) * y \implies e * x = e * y \implies x = y$.
3. Se $M$ è finito e $a$ è cancellabile, la traslazione sinistra $\lambda_a: M \to M$ è iniettiva. Per il Teorema~\ref{teo:teo_finiti_uguali}, un'applicazione iniettiva tra insiemi finiti della stessa cardinalità è anche \textbf{suriettiva}. Poiché $e \in M$, deve esistere un elemento $b \in M$ tale che $\lambda_a(b) = e$, cioè $a * b = e$. Dunque $a$ ammette inverso destro $b$. Analogamente la traslazione destra fornisce un inverso sinistro, da cui $a \in \mathcal{U}(M)$.
\end{proof}

\section{Semigruppi, Monoidi e Gruppi}

\begin{definizione}{Gerarchia delle Strutture ad una Operazione}{gerarchia_gruppi}
Una struttura algebrica $(G, *)$ si dice:
\begin{enumerate}
    \item \textbf{Semigruppo:} se $*$ è associativa.
    \item \textbf{Monoide:} se è un semigruppo dotato di elemento neutro.
    L'insieme degli elementi invertibili di un monoide $M$ si indica con $\mathcal{U}(M)$ e forma un gruppo rispetto a $*$.
    \item \textbf{Gruppo:} se è un monoide in cui \textbf{ogni elemento è simmetrizzabile}:
    \[
        \forall a \in G, \quad \exists a^{-1} \in G : a * a^{-1} = a^{-1} * a = e
    \]
    Se l'operazione è commutativa, $G$ è detto \textbf{gruppo abeliano} (o commutativo).
\end{enumerate}
\end{definizione}

\subsection{Sottogruppi}
Sia $(G, \cdot)$ un gruppo. Un sottoinsieme non vuoto $H \sseq G$ si dice \textbf{sottogruppo} di $G$ (e si scrive $H \le G$) se $H$ è un gruppo rispetto all'operazione indotta.

\begin{teorema}{Criterio di Caratterizzazione dei Sottogruppi}{criterio_sottogruppi}
Sia $G$ un gruppo e sia $H \sseq G$ un sottoinsieme non vuoto ($H \neq \vuoto$). Allora:
\[
    H \le G \quad \Shiff \quad \forall a, b \in H, \quad a b^{-1} \in H
\]
Se $H$ è un sottoinsieme \textbf{finito}, la condizione si semplifica nella sola chiusura:
\[
    H \text{ finito non vuoto} \Implica (H \le G \iff \forall a, b \in H, ab \in H)
\]
\end{teorema}

\subsection{Gruppi Ciclici e Teorema di Lagrange}
Dato un elemento $a \in G$, il \textbf{sottogruppo ciclico generato} da $a$ è:
\[
    \gen{a} = \{a^k \mid k \in \Z\}
\]
L'\textbf{ordine} di $a$, denotato con $\operatorname{ord}(a)$, è la cardinalità $|\gen{a}|$. Se $\gen{a}$ è finito, $\operatorname{ord}(a)$ coincide con il minimo intero positivo $k \ge 1$ tale che $a^k = e$.

\begin{teorema}{Teorema di Lagrange}{lagrange}
Sia $G$ un gruppo finito e sia $H \le G$ un sottogruppo. Allora la cardinalità di $H$ \textbf{divide} la cardinalità di $G$:
\[
    |H| \text{ divide } |G|
\]
Il quoziente $[G : H] \coloneqq |G| / |H|$ è detto \textbf{indice} di $H$ in $G$ e conta il numero delle classi laterali $gH = \{gh \mid h \in H\}$.
\end{teorema}

\begin{proposizione}{Corollari Fondamentali del Teorema di Lagrange}{corollari_lagrange}
Sia $G$ un gruppo finito di ordine $n = |G|$.
\begin{enumerate}
    \item Per ogni elemento $a \in G$, l'ordine dell'elemento $\operatorname{ord}(a)$ divide $|G|$.
    \item Per ogni $a \in G$, si ha $a^{|G|} = e$.
    \item Ogni gruppo di ordine $p$ con $p$ \textbf{numero primo} è un \textbf{gruppo ciclico} (e quindi abeliano), privo di sottogruppi propri non banali.
\end{enumerate}
\end{proposizione}

\subsection{Omomorfismi di Gruppi}

\begin{definizione}{Omomorfismo di Gruppi}{omom_gruppi}
Dati due gruppi $(G, \cdot)$ e $(G', *)$, un'applicazione $\varphi: G \to G'$ è detta \textbf{omomorfismo di gruppi} se preserva l'operazione:
\[
    \forall a, b \in G, \quad \varphi(ab) = \varphi(a) * \varphi(b)
\]
\end{definizione}

Per ogni omomorfismo di gruppi si ha $\varphi(e_G) = e_{G'}$ e $\varphi(a^{-1}) = (\varphi(a))^{-1}$.
\begin{itemize}
    \item Il \textbf{nucleo} di $\varphi$ è il sottogruppo $\Ker(\varphi) = \{x \in G \mid \varphi(x) = e_{G'}\} \le G$.
    \item L'\textbf{immagine} di $\varphi$ è il sottogruppo $\im(\varphi) = \{\varphi(x) \mid x \in G\} \le G'$.
\end{itemize}

\begin{teorema}{Caratterizzazione dell'Iniettività tramite il Nucleo}{nucleo_iniettivita}
Un omomorfismo di gruppi $\varphi: G \to G'$ è \textbf{iniettivo} se e solo se il suo nucleo è banale:
\[
    \varphi \text{ è iniettivo} \quad \Shiff \quad \Ker(\varphi) = \{e_G\}
\]
\end{teorema}
\begin{proof}
$(\implica)$ Se $\varphi$ è iniettivo, poiché $\varphi(e_G) = e_{G'}$, nessun altro elemento può essere mappato in $e_{G'}$, quindi $\Ker(\varphi) = \{e_G\}$.
$(\impliedby)$ Sia $\Ker(\varphi) = \{e_G\}$. Siano $a, b \in G$ tali che $\varphi(a) = \varphi(b)$.
Moltiplicando per $\varphi(b)^{-1}$: $\varphi(a) * \varphi(b)^{-1} = e_{G'} \implies \varphi(ab^{-1}) = e_{G'}$.
Ne segue che $ab^{-1} \in \Ker(\varphi) = \{e_G\}$, da cui $ab^{-1} = e_G$, cioè $a = b$. Dunque $\varphi$ è iniettivo.
\end{proof}

\section{Anelli, Domini d'Integrità e Campi}

\begin{definizione}{Anello}{anello}
Una struttura algebrica $(R, +, \cdot)$ dotata di due operazioni binarie interne (dette \emph{addizione} e \emph{moltiplicazione}) si dice \textbf{anello} se:
\begin{enumerate}
    \item $(R, +)$ è un \textbf{gruppo abeliano} (con elemento neutro denotato con $0$ e opposto $-a$);
    \item $(R, \cdot)$ è un \textbf{semigruppo} (la moltiplicazione è associativa);
    \item La moltiplicazione è \textbf{distributiva} rispetto all'addizione (a sinistra e a destra):
    \[
        \forall a, b, c \in R, \quad a \cdot (b + c) = a \cdot b + a \cdot c, \qquad (a + b) \cdot c = a \cdot c + b \cdot c
    \]
\end{enumerate}
Un anello $R$ si dice:
\begin{itemize}
    \item \textbf{Commutativo:} se la moltiplicazione è commutativa ($ab = ba$).
    \item \textbf{Unitario:} se possiede un elemento neutro moltiplicativo (denotato con $1 \neq 0$).
\end{itemize}
\end{definizione}

In ogni anello l'elemento $0$ è moltiplicativamente \textbf{assorbente}: per ogni $a \in R$, $a \cdot 0 = 0 \cdot a = 0$. Infatti $a \cdot 0 = a \cdot (0 + 0) = a \cdot 0 + a \cdot 0 \implies a \cdot 0 = 0$ per cancellabilità nel gruppo additivo.

\subsection{Divisori dello Zero e Domini d'Integrità}

\begin{definizione}{Divisori dello Zero e Dominio d'Integrità}{div_zero}
Sia $R$ un anello commutativo unitario.
\begin{itemize}
    \item Un elemento non nullo $a \in R \smin \{0\}$ si dice \textbf{divisore dello zero} se esiste un elemento non nullo $b \in R \smin \{0\}$ tale che:
    \[
        a \cdot b = 0
    \]
    \item Un elemento $a \in R$ è \textbf{cancellabile} (rispetto al prodotto) se e solo se \textbf{non è un divisore dello zero}.
    \item Un anello commutativo unitario $R$ privo di divisori dello zero (cioè in cui $ab = 0 \implica a = 0 \vee b = 0$) è detto \textbf{dominio d'integrità} (o anello integro).
\end{itemize}
\end{definizione}

\begin{definizione}{Corpi e Campi}{campi}
Un anello unitario $(K, +, \cdot)$ si dice:
\begin{itemize}
    \item \textbf{Corpo} (o anello con divisione) se ogni elemento non nullo ammette inverso moltiplicativo:
    \[
        \mathcal{U}(K) = K \smin \{0\} = K^*
    \]
    \item \textbf{Campo} (in inglese \emph{field}) se è un corpo commutativo.
\end{itemize}
\end{definizione}

Ogni campo è un dominio d'integrità. Infatti, se $ab = 0$ con $a \neq 0$, moltiplicando per $a^{-1}$ si ha $a^{-1}(ab) = a^{-1} \cdot 0 \implies (a^{-1}a)b = 0 \implies 1 \cdot b = 0 \implies b = 0$.

\begin{teorema}{Teorema di Wedderburn per Domini Finiti}{domini_finiti}
Ogni dominio d'integrità \textbf{finito} è un \textbf{campo}.
\end{teorema}
\begin{proof}
Sia $R$ un dominio d'integrità finito e sia $a \in R$ con $a \neq 0$.
Consideriamo la traslazione moltiplicativa $\lambda_a: R \to R$ definita da $\lambda_a(x) = ax$.
Mostriamo che $\lambda_a$ è iniettiva: se $ax_1 = ax_2$, allora $a(x_1 - x_2) = 0$. Poiché $R$ è un dominio e $a \neq 0$, deve essere $x_1 - x_2 = 0$, da cui $x_1 = x_2$.
Essendo $R$ finito, per il Teorema~\ref{teo:teo_finiti_uguali} ogni applicazione iniettiva $\lambda_a: R \to R$ è anche suriettiva.
Poiché $1 \in R$, deve esistere un elemento $b \in R$ tale che $\lambda_a(b) = 1$, cioè $ab = 1$. Dunque $a$ è invertibile.
Essendo $a$ un generico elemento non nullo, ogni elemento non nullo è invertibile, dunque $R$ è un campo.
\end{proof}

\subsection{Sottoanelli ed Ideali}
Sia $R$ un anello commutativo unitario.
\begin{itemize}
    \item Un sottoinsieme $S \sseq R$ è un \textbf{sottoanello} se $(S, +) \le (R, +)$ e $S$ è chiuso rispetto al prodotto ($s_1 s_2 \in S$) e contiene $1_R$.
    \item Un sottoinsieme $I \sseq R$ è un \textbf{ideale} (scritto $I \trianglelefteq R$) se:
    \begin{enumerate}
        \item $(I, +)$ è un sottogruppo di $(R, +)$ ($0 \in I$ e $\forall x, y \in I, x - y \in I$);
        \item Proprietà assorbente: per ogni $r \in R$ e ogni $x \in I$, si ha $rx \in I$.
    \end{enumerate}
\end{itemize}

Dato un ideale $I \trianglelefteq R$, l'insieme quoziente $R/I = \{a + I \mid a \in R\}$ eredita una struttura di anello con le operazioni $(a + I) + (b + I) = (a + b) + I$ e $(a + I)(b + I) = ab + I$.
Se $\varphi: R \to R'$ è un omomorfismo di anelli, il suo nucleo $\Ker(\varphi) = \{x \in R \mid \varphi(x) = 0\}$ è un ideale di $R$, e vale il \textbf{Primo Teorema di Isomorfismo per Anelli}:
\[
    R/\Ker(\varphi) \cong \im(\varphi)
\]

\chapter{Anelli di Polinomi e Loro Aritmetica}
\label{chap:polinomi}

In questo capitolo studiamo l'anello dei polinomi $K[x]$ a coefficienti in un campo $K$. L'aritmetica dei polinomi presenta profonde analogie con l'aritmetica di $\Z$: divisione euclidea, algoritmo di Bézout, elementi irriducibili e costruzione di campi quoziente.

\section{Costruzione dell'Anello dei Polinomi}

Sia $A$ un anello commutativo unitario.

\begin{definizione}{Polinomio a Coefficienti in $A$}{polinomio}
Un \textbf{polinomio} nell'indeterminata $x$ a coefficienti in $A$ è una successione di elementi di $A$, $f = (a_0, a_1, a_2, \dots, a_n, \dots)$, definitivamente nulla (ovvero tale che esiste $N \in \N$ per cui $a_k = 0$ per ogni $k > N$).

Denotando con $x = (0, 1, 0, 0, \dots)$ l'indeterminata, ogni polinomio non nullo si scrive in modo unico nella forma canonica:
\[
    f(x) = a_0 + a_1 x + a_2 x^2 + \dots + a_n x^n = \sum_{i=0}^n a_i x^i \qquad (a_n \neq 0)
\]
L'insieme di tutti i polinomi nell'indeterminata $x$ a coefficienti in $A$ si denota con $A[x]$.
\end{definizione}

\begin{definizione}{Grado e Coefficiente Direttivo}{grado}
Sia $f(x) = \sum_{i=0}^n a_i x^i \in A[x]$ un polinomio non nullo ($f \neq 0$):
\begin{itemize}
    \item Il massimo indice $n \in \N$ per cui $a_n \neq 0$ è detto \textbf{grado} del polinomio e si denota con $\degpoly(f) = n$.
    \item $a_n \in A \smin \{0\}$ è detto \textbf{coefficiente direttivo} (o principale).
    \item Se $a_n = 1$, il polinomio si dice \textbf{monico}.
    \item Il polinomio nullo $0$ non ha coefficienti non nulli: per convenzione formale si pone $\degpoly(0) = -\infty$.
\end{itemize}
\end{definizione}

\begin{teorema}{Proprietà del Grado e Integrità di $A[x]$}{grado_integrita}
Siano $f, g \in A[x]$:
\begin{enumerate}
    \item $\degpoly(f + g) \le \max(\degpoly(f), \degpoly(g))$ (l'uguaglianza vale se $\degpoly(f) \neq \degpoly(g)$).
    \item $\degpoly(f \cdot g) \le \degpoly(f) + \degpoly(g)$.
    \item Se l'anello dei coefficienti $A$ è un \textbf{dominio d'integrità}, allora per ogni $f, g \in A[x] \smin \{0\}$:
    \[
        \degpoly(f \cdot g) = \degpoly(f) + \degpoly(g)
    \]
    Di conseguenza, $A[x]$ è anch'esso un \textbf{dominio d'integrità}.
    \item Inoltre, gli elementi invertibili di $A[x]$ coincidono con le sole costanti invertibili di $A$:
    \[
        \mathcal{U}(A[x]) = \mathcal{U}(A)
    \]
    In particolare, se $K$ è un campo, $\mathcal{U}(K[x]) = K^* = K \smin \{0\}$ (tutti i polinomi di grado 0).
\end{enumerate}
\end{teorema}

\section{Divisione Euclidea e Algoritmo di Bézout in $K[x]$}

Sia $K$ un campo (ad esempio $\Q, \R, \C$ oppure $\Z_p$ con $p$ primo).

\begin{teorema}{Divisione con Resto in $K[x]$}{div_polinomi}
Siano $f(x), g(x) \in K[x]$ con $g(x) \neq 0$. Esiste un'\textbf{unica} coppia di polinomi $q(x), r(x) \in K[x]$ (detti rispettivamente \textbf{quoziente} e \textbf{resto}) tali che:
\[
    f(x) = g(x) q(x) + r(x) \qquad \text{con } \degpoly(r) < \degpoly(g)
\]
\end{teorema}

\begin{definizione}{Massimo Comune Divisore tra Polinomi}{mcd_polinomi}
Dati $f, g \in K[x]$ non entrambi nulli, un polinomio \textbf{monico} $d(x) \in K[x]$ è detto \textbf{massimo comune divisore} di $f$ e $g$ (denotato con $\mcd(f, g)$) se:
\begin{enumerate}
    \item $d(x) \mid f(x)$ e $d(x) \mid g(x)$;
    \item Per ogni $c(x) \in K[x]$, se $c(x) \mid f(x)$ e $c(x) \mid g(x)$, allora $c(x) \mid d(x)$.
\end{enumerate}
Per il \textbf{Teorema di Bézout per polinomi}, esistono $u(x), v(x) \in K[x]$ tali che:
\[
    d(x) = f(x)u(x) + g(x)v(x)
\]
\end{definizione}

Il calcolo di $\mcd(f, g)$ e dei polinomi di Bézout $u, v$ avviene applicando l'algoritmo euclideo formale per polinomi.

\section{Radici di un Polinomio e Teorema di Ruffini}

Sia $f(x) = \sum_{i=0}^n a_i x^i \in K[x]$. L'\textbf{omomorfismo di valutazione} associato a un elemento $c \in K$ è:
\[
    \operatorname{ev}_c: K[x] \longrightarrow K, \qquad f(x) \longmapsto f(c) = \sum_{i=0}^n a_i c^i
\]
Un elemento $c \in K$ si dice \textbf{radice} (o zero) di $f(x)$ se $f(c) = 0$.

\begin{teorema}{Teorema del Resto e Teorema di Ruffini}{ruffini}
Siano $f(x) \in K[x]$ e $c \in K$.
\begin{enumerate}
    \item \textbf{Teorema del Resto:} Il resto della divisione di $f(x)$ per il binomio $(x - c)$ è la costante $f(c) \in K$:
    \[
        f(x) = (x - c)q(x) + f(c)
    \]
    \item \textbf{Teorema di Ruffini:} Un elemento $c \in K$ è radice di $f(x)$ se e solo se il binomio $(x - c)$ divide $f(x)$:
    \[
        f(c) = 0 \quad \Shiff \quad (x - c) \mid f(x) \text{ in } K[x]
    \]
\end{enumerate}
\end{teorema}
\begin{proof}
Dividendo $f(x)$ per $(x - c)$ tramite il Teorema~\ref{teo:div_polinomi}, otteniamo $f(x) = (x - c)q(x) + r(x)$ con $\degpoly(r) < \degpoly(x - c) = 1$.
Dunque $r(x) = r_0$ è una costante appartenente a $K$.
Valutando l'uguaglianza in $x = c$:
\[
    f(c) = (c - c)q(c) + r_0 = 0 \cdot q(c) + r_0 = r_0
\]
Ne segue che $f(x) = (x - c)q(x) + f(c)$. Da ciò discende immediatamente che $(x - c) \mid f(x) \iff f(c) = 0$.
\end{proof}

\begin{definizione}{Molteplicità di una Radice}{molteplicita}
Sia $c \in K$ una radice di $f(x) \in K[x] \smin \{0\}$. La \textbf{molteplicità} di $c$ è il massimo intero positivo $m \ge 1$ tale che $(x - c)^m$ divide $f(x)$ in $K[x]$:
\[
    f(x) = (x - c)^m g(x) \qquad \text{con } g(c) \neq 0
\]
\end{definizione}

\begin{teorema}{Numero Massimo di Radici}{max_radici}
Un polinomio non nullo $f(x) \in K[x]$ di grado $n = \degpoly(f)$ a coefficienti in un campo $K$ ammette \textbf{al più $n$ radici} in $K$, contate con la loro molteplicità.
\end{teorema}

\section{Polinomi Irriducibili e Fattorizzazione}

\begin{definizione}{Polinomio Irriducibile}{irriducibile}
Un polinomio non costante $f(x) \in K[x]$ ($\degpoly(f) \ge 1$) si dice \textbf{irriducibile} in $K[x]$ se non può essere espresso come prodotto di due polinomi di grado strettamente positivo:
\[
    f(x) = g(x) h(x) \Implica (\degpoly(g) = 0 \vee \degpoly(h) = 0)
\]
Un polinomio non costante che non è irriducibile si dice \textbf{riducibile}.
\end{definizione}

\begin{proposizione}{Irriducibilità per Gradi 2 e 3}{criterio_grado_2_3}
Sia $f(x) \in K[x]$ un polinomio di grado $2$ oppure $3$. Allora:
\[
    f(x) \text{ è irriducibile in } K[x] \quad \Shiff \quad f(x) \text{ non ammette alcuna radice in } K
\]
\end{proposizione}
\begin{proof}
Se $f$ è riducibile, allora $f(x) = g(x)h(x)$ con $\degpoly(g), \degpoly(h) \ge 1$.
Poiché $\degpoly(g) + \degpoly(h) \in \{2, 3\}$, almeno uno dei due fattori deve avere grado $1$, cioè $g(x) = ax + b$ con $a \neq 0$.
Allora $c = -a^{-1}b \in K$ è una radice di $g(x)$ e quindi di $f(x)$.
Per contrapposizione, se $f$ non ha radici in $K$, non può ammettere fattori lineari, dunque è irriducibile.
\end{proof}

\begin{avvisoCutolo}
    \textbf{Attenzione: Il criterio delle radici NON vale per grado $\ge 4$!}
    Un polinomio di grado 4 può non avere alcuna radice ed essere comunque riducibile, spezzandosi nel prodotto di due polinomi irriducibili di secondo grado!
    \emph{Esempio classico in $\R[x]$:} $f(x) = x^4 + 2x^2 + 1 = (x^2 + 1)^2$. Non ha radici in $\R$, ma è riducibile!
\end{avvisoCutolo}

\subsection{Criteri di Irriducibilità nei Vari Campi}
\begin{enumerate}
    \item \textbf{In $\C[x]$:} per il \emph{Teorema Fondamentale dell'Algebra}, gli unici polinomi irriducibili sono quelli di \textbf{primo grado} ($ax + b$).
    \item \textbf{In $\R[x]$:} gli unici polinomi irriducibili sono quelli di \textbf{primo grado} e quelli di \textbf{secondo grado con discriminante negativo} ($\Delta = b^2 - 4ac < 0$).
    \item \textbf{In $\Q[x]$:}
    \begin{itemize}
        \item \textbf{Criterio delle radici razionali:} se $f(x) = a_n x^n + \dots + a_0 \in \Z[x]$ ha radice razionale $p/q$ con $\mcd(p, q) = 1$, allora $p \mid a_0$ e $q \mid a_n$.
        \item \textbf{Criterio di Eisenstein:} se esiste un primo $p$ tale che $p \nmid a_n$, $p \mid a_i$ per ogni $i \in \{0, \dots, n-1\}$ e $p^2 \nmid a_0$, allora $f(x)$ è irriducibile in $\Q[x]$.
        \item \textbf{Riduzione modulo $p$:} se riducendo i coefficienti modulo un primo $p$ il grado si conserva e il polinomio è irriducibile in $\Z_p[x]$, allora è irriducibile in $\Q[x]$.
    \end{itemize}
    \item \textbf{In $\Z_p[x]$ ($p$ primo):} l'insieme $\Z_p = \{0, 1, \dots, p-1\}$ è finito. Per testare l'irriducibilità di un polinomio di secondo o terzo grado, si valuta il polinomio in tutti i $p$ elementi: se $f(c) \not\equiv 0 \pmod p$ per ogni $c \in \Z_p$, il polinomio è irriducibile.
\end{enumerate}

\section{L'Anello Quoziente $K[x]/(f(x))$ e Costruzione di Campi}

Dato un polinomio non costante $f(x) \in K[x]$, consideriamo l'ideale principale generato da $f$:
\[
    (f(x)) = \{f(x) \cdot g(x) \mid g(x) \in K[x]\}
\]
L'anello quoziente $K[x]/(f(x))$ è costituito dalle classi di equivalenza dei resti della divisione per $f(x)$:
\[
    K[x]/(f(x)) = \{[r(x)] \mid r(x) \in K[x], \degpoly(r) < \degpoly(f)\}
\]
Se $\degpoly(f) = n$, come spazio vettoriale su $K$ ha dimensione $n$, con base $\{[1], [x], [x^2], \dots, [x^{n-1}]\}$.

\begin{teorema}{Caratterizzazione dei Campi Quoziente di Polinomi}{campo_quoziente}
Sia $K$ un campo e sia $f(x) \in K[x]$ un polinomio non costante.
L'anello quoziente $K[x]/(f(x))$ è un \textbf{campo} se e solo se $f(x)$ è un polinomio \textbf{irriducibile} in $K[x]$:
\[
    K[x]/(f(x)) \text{ è un campo} \quad \Shiff \quad f(x) \text{ è irriducibile in } K[x]
\]
\end{teorema}
\begin{proof}
$(\implica)$ Se $f$ fosse riducibile, $f(x) = g(x)h(x)$ con $1 \le \degpoly(g), \degpoly(h) < \degpoly(f)$.
Passando alle classi nel quoziente: $[g(x)] \cdot [h(x)] = [f(x)] = [0]$.
Poiché $\degpoly(g), \degpoly(h) < \degpoly(f)$, si ha $[g] \neq [0]$ e $[h] \neq [0]$.
Dunque $[g]$ e $[h]$ sarebbero divisori dello zero, impedendo a $K[x]/(f(x))$ di essere un dominio, e tanto meno un campo.
$(\impliedby)$ Sia $f$ irriducibile e sia $[g(x)] \in K[x]/(f(x))$ con $[g] \neq [0]$.
Allora $f(x) \nmid g(x)$. Poiché $f(x)$ è irriducibile, i suoi unici divisori sono le costanti e i suoi associati, dunque $\mcd(f(x), g(x)) = 1$.
Per il Teorema di Bézout per polinomi, esistono $u(x), v(x) \in K[x]$ tali che $f(x)u(x) + g(x)v(x) = 1$.
Passando al quoziente modulo $f(x)$:
\[
    [g(x)] \cdot [v(x)] = [1]
\]
Dunque $[g(x)]$ è invertibile con inverso $[v(x)]$. Ne segue che ogni elemento non nullo è invertibile, quindi $K[x]/(f(x))$ è un campo.
\end{proof}


% ------------------------------------------------------------------------------
% PARTE V: TEORIA DEI GRAFI
% ------------------------------------------------------------------------------
\part{Teoria dei Grafi}
\chapter{Elementi di Teoria dei Grafi}
\label{chap:grafi}

La teoria dei grafi fornisce modelli matematici per descrivere relazioni binarie simmetriche tra oggetti discreti. In questo capitolo presentiamo la teoria fondamentale: gradi dei vertici, lemma delle strette di mano, alberi, grafi bipartiti, cammini euleriani e isomorfismi.

\section{Definizioni Fondamentali}

\begin{definizione}{Grafo Semplice non Orientato}{grafo}
Un \textbf{grafo semplice non orientato} è una coppia ordinata $G = (V, E)$, dove:
\begin{itemize}
    \item $V$ è un insieme non vuoto, i cui elementi sono detti \textbf{vertici} (o nodi);
    \item $E \sseq \P_2(V) = \{\{u, v\} \sseq V \mid u \neq v\}$ è un insieme di sottoinsiemi a due elementi di $V$, detti \textbf{spigoli} (o archi, \emph{edges}).
\end{itemize}
Se $\{u, v\} \in E$, si dice che i vertici $u$ e $v$ sono \textbf{adiacenti} (o vicini) e che lo spigolo $\{u, v\}$ è \textbf{incidente} sui vertici $u$ e $v$.
\end{definizione}

Un grafo si dice \textbf{finito} se l'insieme dei vertici $V$ è finito.
L'\textbf{ordine} del grafo è il numero dei vertici $n = |V|$; la \textbf{dimensione} è il numero degli spigoli $m = |E|$.

\subsection{Grado di un Vertice e Lemma delle Strette di Mano}

\begin{definizione}{Grado di un Vertice}{grado_vertice}
Sia $G = (V, E)$ un grafo e sia $v \in V$. Il \textbf{grado} (o valenza) di $v$, denotato con $\deg(v)$ o $d(v)$, è il numero di spigoli incidenti su $v$:
\[
    \deg(v) = |\{e \in E \mid v \in e\}|
\]
Un vertice con $\deg(v) = 0$ si dice \textbf{isolato}; un vertice con $\deg(v) = 1$ si dice \textbf{pendente} (o foglia).
\end{definizione}

\begin{teorema}{Lemma delle Strette di Mano (Handshaking Lemma)}{handshaking}
In ogni grafo finito $G = (V, E)$, la somma dei gradi di tutti i vertici è uguale al \textbf{doppio del numero degli spigoli}:
\[
    \sum_{v \in V} \deg(v) = 2 |E|
\]
\end{teorema}
\begin{proof}
Ogni spigolo $e = \{u, v\} \in E$ possiede esattamente due estremi, $u$ e $v$. Di conseguenza, nel conteggio della somma $\sum_{v \in V} \deg(v)$, ogni spigolo viene conteggiato esattamente due volte (una per l'estremo $u$ e una per l'estremo $v$). Dunque la somma coincide con $2|E|$.
\end{proof}

\begin{proposizione}{Parità dei Vertici di Grado Dispari}{vertici_dispari}
In qualunque grafo finito, il numero di vertici aventi grado dispari è sempre \textbf{pari}.
\end{proposizione}
\begin{proof}
Partizioniamo l'insieme dei vertici $V$ in due sottoinsiemi disgiunti: $V_{\text{pari}} = \{v \in V \mid \deg(v) \text{ è pari}\}$ e $V_{\text{dispari}} = \{v \in V \mid \deg(v) \text{ è dispari}\}$.
Per il Lemma delle Strette di Mano:
\[
    \sum_{v \in V_{\text{pari}}} \deg(v) + \sum_{v \in V_{\text{dispari}}} \deg(v) = \sum_{v \in V} \deg(v) = 2|E| \quad (\text{numero pari})
\]
La prima sommatoria è somma di numeri pari, dunque è pari. Ne segue che anche la seconda sommatoria $\sum_{v \in V_{\text{dispari}}} \deg(v)$ deve essere un numero pari.
Poiché una somma di numeri dispari è pari se e solo se il numero degli addendi è pari, deduciamo che la cardinalità $|V_{\text{dispari}}|$ è necessariamente un numero pari.
\end{proof}

\section{Connessione, Cammini e Cicli}

\begin{definizione}{Cammini e Cicli}{cammini_cicli}
Sia $G = (V, E)$ un grafo.
\begin{itemize}
    \item Un \textbf{cammino} di lunghezza $k \ge 0$ tra $v_0$ e $v_k$ è una successione di vertici $(v_0, v_1, \dots, v_k)$ tale che $\{v_{i-1}, v_i\} \in E$ per ogni $i \in \{1, \dots, k\}$.
    \item Il cammino si dice \textbf{semplice} se non ripete alcun vertice (tutti i $v_i$ sono distinti).
    \item Un cammino si dice \textbf{chiuso} se $v_0 = v_k$.
    \item Un \textbf{ciclo} di lunghezza $k \ge 3$ è un cammino chiuso $(v_0, v_1, \dots, v_{k-1}, v_0)$ in cui i vertici $v_0, v_1, \dots, v_{k-1}$ sono a due a due distinti.
    Un grafo privo di cicli si dice \textbf{aciclico} (o foresta).
\end{itemize}
\end{definizione}

\begin{definizione}{Grafo Connesso e Componenti Connesse}{connessione}
Un grafo $G$ si dice \textbf{connesso} se per ogni coppia di vertici distinti $u, v \in V$ esiste almeno un cammino che li congiunge.
La relazione binaria $\sim$ su $V$ definita da:
\[
    u \sim v \quad \Shiff \quad u = v \;\vee\; \text{esiste un cammino tra } u \text{ e } v
\]
è una \textbf{relazione di equivalenza}. Le sue classi di equivalenza inducono i sottografi connessi massimali di $G$, detti \textbf{componenti connesse}.
\end{definizione}

\section{Alberi e Loro Caratterizzazioni}

\begin{definizione}{Albero}{albero}
Un grafo finito $G = (V, E)$ si dice \textbf{albero} se è \textbf{connesso} e \textbf{privo di cicli} (aciclico).
\end{definizione}

\begin{teorema}{Caratterizzazioni Equivalenti degli Alberi}{teo_alberi}
Sia $G = (V, E)$ un grafo finito con $n = |V| \ge 1$ vertici. Le seguenti affermazioni sono tra loro \textbf{logicamente equivalenti}:
\begin{enumerate}
    \item $G$ è un albero (connesso e aciclico).
    \item Per ogni coppia di vertici distinti $u, v \in V$, esiste un \textbf{unico} cammino semplice che li congiunge.
    \item $G$ è connesso e possiede esattamente $n - 1$ spigoli:
    \[
        |E| = |V| - 1
    \]
    \item $G$ è aciclico e possiede esattamente $n - 1$ spigoli.
    \item $G$ è aciclico, ma aggiungendo un qualunque nuovo spigolo non presente in $E$ si crea esattamente un ciclo.
    \item $G$ è connesso, ma rimuovendo un qualunque spigolo si ottiene un grafo non connesso.
\end{enumerate}
\end{teorema}

\begin{proposizione}{Presenza di Foglie negli Alberi}{foglie_albero}
Ogni albero finito $T = (V, E)$ con almeno due vertici ($|V| \ge 2$) ammette \textbf{almeno due vertici di grado 1} (foglie).
\end{proposizione}

\section{Grafi Bipartiti e Teorema di Kőnig}

\begin{definizione}{Grafo Bipartito}{grafo_bipartito}
Un grafo $G = (V, E)$ si dice \textbf{bipartito} se l'insieme dei suoi vertici $V$ può essere partizionato in due sottoinsiemi disgiunti $V_1$ e $V_2$ ($V = V_1 \cup V_2, V_1 \cap V_2 = \vuoto$) in modo tale che ogni spigolo di $G$ congiunga un vertice di $V_1$ con un vertice di $V_2$:
\[
    \forall e \in E, \quad e \cap V_1 \neq \vuoto \;\wedge\; e \cap V_2 \neq \vuoto
\]
In altre parole, non esiste alcuno spigolo interno a $V_1$ o interno a $V_2$.
\end{definizione}

\begin{teorema}{Criterio di Bipartizione di Kőnig}{konig}
Un grafo $G$ è bipartito \textbf{se e solo se non contiene alcun ciclo di lunghezza dispari}.
\end{teorema}
\begin{proof}
$(\implica)$ Sia $G$ bipartito con partizione $(V_1, V_2)$. Sia $C = (v_0, v_1, \dots, v_k = v_0)$ un ciclo.
Partendo da $v_0 \in V_1$, per definizione di bipartizione si ha $v_1 \in V_2, v_2 \in V_1, v_3 \in V_2$, e in generale i vertici di indice pari stanno in $V_1$ e quelli di indice dispari in $V_2$.
Poiché il ciclo si chiude su $v_k = v_0 \in V_1$, l'indice $k$ deve essere necessariamente un numero \textbf{pari}. Dunque ogni ciclo ha lunghezza pari.
$(\impliedby)$ Si assume $G$ privo di cicli dispari (possiamo supporre $G$ connesso). Fissato un vertice radice $v_0$, si definisce la distanza $d(v_0, x)$ come la lunghezza minima di un cammino tra $v_0$ e $x$. Ponendo $V_1 = \{x \in V \mid d(v_0, x) \text{ è pari}\}$ e $V_2 = \{x \in V \mid d(v_0, x) \text{ è dispari}\}$, l'assenza di cicli dispari garantisce che nessun vertice di $V_1$ sia adiacente a un altro vertice di $V_1$, né in $V_2$. Dunque $G$ è bipartito.
\end{proof}

\section{Cammini e Circuiti Euleriani}

Il problema dei ponti di Königsberg diede origine alla teoria dei grafi nel 1736 ad opera di Leonhard Euler.

\begin{definizione}{Cammino e Circuito Euleriano}{euleriano}
Sia $G = (V, E)$ un grafo finito:
\begin{itemize}
    \item Un \textbf{cammino euleriano} è un cammino che attraversa \textbf{ogni spigolo di $G$ esattamente una volta}.
    \item Un \textbf{circuito euleriano} è un cammino euleriano chiuso (che inizia e termina nello stesso vertice).
    Un grafo che ammette un circuito euleriano è detto \textbf{grafo euleriano}.
\end{itemize}
\end{definizione}

\begin{teorema}{Teorema di Eulero per i Grafi Euleriani}{teo_eulero_grafi}
Sia $G = (V, E)$ un grafo connesso (a meno di vertici isolati):
\begin{enumerate}
    \item $G$ ammette un \textbf{circuito euleriano} se e solo se \textbf{tutti i vertici hanno grado pari}:
    \[
        \forall v \in V, \quad \deg(v) \text{ è pari}
    \]
    \item $G$ ammette un \textbf{cammino euleriano aperto} (con estremi distinti) se e solo se ammette \textbf{esattamente due vertici di grado dispari}. In tal caso, il cammino deve necessariamente iniziare in uno dei due vertici dispari e terminare nell'altro.
\end{enumerate}
\end{teorema}

\section{Isomorfismo di Grafi ed Invarianti}

\begin{definizione}{Isomorfismo di Grafi}{isomorfismo_grafi}
Due grafi $G_1 = (V_1, E_1)$ e $G_2 = (V_2, E_2)$ si dicono \textbf{isomorfi} (e si scrive $G_1 \cong G_2$) se esiste una \textbf{biiezione} tra i rispettivi insiemi di vertici:
\[
    f: V_1 \longrightarrow V_2
\]
tale che due vertici sono adiacenti in $G_1$ se e solo se le loro immagini sono adiacenti in $G_2$:
\[
    \forall u, v \in V_1, \quad \{u, v\} \in E_1 \quad \Shiff \quad \{f(u), f(v)\} \in E_2
\]
L'applicazione $f$ è detta \textbf{isomorfismo di grafi}.
\end{definizione}

\subsection{Invarianti di Isomorfismo}
Dimostrare che due grafi sono isomorfi richiede di esibire esplicitamente la biiezione $f$.
Per dimostrare che due grafi \textbf{non sono isomorfi}, la strategia più rapida consiste nell'individuare una proprietà strutturale (detta \textbf{invariante di isomorfismo}) posseduta da uno ma non dall'altro.

\begin{proposizione}{Principali Invarianti di Isomorfismo}{invarianti}
Se due grafi $G_1$ e $G_2$ sono isomorfi ($G_1 \cong G_2$), allora devono condividere:
\begin{enumerate}
    \item Lo stesso numero di vertici: $|V_1| = |V_2|$;
    \item Lo stesso numero di spigoli: $|E_1| = |E_2|$;
    \item La stessa \textbf{sequenza dei gradi} ordinata (ad es. $(3, 3, 2, 2, 2)$);
    \item Lo stesso numero di componenti connesse;
    \item Lo stesso numero di cicli di ciascuna lunghezza prefissata $k$ (in particolare, triangoli di lunghezza 3 e quadrati di lunghezza 4);
    \item La stessa lunghezza del ciclo minimo (\emph{girth});
    \item La proprietà di essere o meno bipartito (o euleriano).
\end{enumerate}
\end{proposizione}


% ------------------------------------------------------------------------------
% PARTE VI: PREPARAZIONE ESAME ED ESERCIZIARIO RISOLTO
% ------------------------------------------------------------------------------
\part{Vademecum d'Esame ed Eserciziario Svolto}
\chapter{Vademecum Strategico per l'Esame: Scritto e Orale}
\label{chap:guida_esame}

L'esame di Algebra con il Prof. Giovanni Cutolo è considerato uno degli scogli più impegnativi del primo anno di Informatica. Questa difficoltà non deriva da calcoli lunghi o noiosi, bensì dall'\textbf{assoluta intransigenza del docente rispetto al rigore linguistico e logico}. In questo capitolo raccogliamo i consigli strategici, l'analisi della prova scritta e il censimento completo delle domande d'orale più insidiose.

\section{La Prova Scritta: Struttura e Gestione dei Tempi}

La prova scritta (della durata di 2 ore nel Nuovo Ordinamento) presenta una struttura standardizzata:

\begin{enumerate}
    \item \textbf{Sezione 1: Quesiti a Risposta Aperta Sintetica (Quesiti Teorici)}
    \begin{itemize}
        \item Da 4 a 6 domande brevi mirate a verificare la conoscenza puntuale delle definizioni assiomatiche e dei teoremi fondamentali (ad es. enunciare il Lemma di Euclide, definire un reticolo booleano, scrivere la formula di $\varphi(m)$, indicare gli invertibili di $\Z_m$).
        \item \emph{Consiglio tattico:} Rispondere con precisione chirurgica. Se la domanda chiede una definizione, scrivere la formula con tutti i quantificatori ($\forall, \exists$) al posto giusto. Non aggiungere divagazioni non richieste che potrebbero contenere imprecisioni.
    \end{itemize}
    \item \textbf{Sezione 2: Esercizio sulle Relazioni d'Ordine, Poset e Reticoli}
    \begin{itemize}
        \item Assegnato un insieme $S$ (ad es. l'insieme dei divisori di un numero, o un sottoinsieme di $\P(A)$) con una relazione $\le$, viene richiesto di: verificare se è un ordine, tracciare il diagramma di Hasse, determinare minimali, massimali, minimo, massimo, calcolare $\inf$ e $\sup$ di una coppia, verificare se è un reticolo e studiarne distributività e complementi.
    \end{itemize}
    \item \textbf{Sezione 3: Esercizio di Aritmetica Modulare}
    \begin{itemize}
        \item Risoluzione di un'equazione congruenziale lineare $ax \equiv b \pmod m$ o di un sistema con il Teorema Cinese del Resto.
        \item Calcolo di una potenza modulare elevata ($a^N \pmod m$) mediante riduzione dell'esponente con il Teorema di Eulero-Fermat o di Fermat.
    \end{itemize}
    \item \textbf{Sezione 4: Strutture Algebriche / Polinomi / Grafi}
    \begin{itemize}
        \item Studio di un'operazione definita su un insieme $S$, verifica degli assiomi di gruppo o anello, calcolo degli elementi invertibili e cancellabili.
        \item Decomposizione in irriducibili di un polinomio in $\Z_p[x]$ e studio del campo quoziente $\Z_p[x]/(f(x))$.
        \item Verifica di proprietà di un grafo (connessione, alberi, bipartizione con Kőnig, eulerianità).
    \end{itemize}
\end{enumerate}

\section{Le 10 Domande ``Killer'' dell'Orale con Cutolo}

Dall'analisi delle esperienze d'esame e delle trascrizioni dei colloqui, emergono dieci domande ricorrenti su cui il docente saggia la vera profondità di comprensione dello studente.

\subsection{1. Perché i Naturali contengono lo zero? ($\N$ e Assiomi di Peano)}
\textbf{La trappola:} Lo studente risponde ``perché $0$ è il primo numero''.
\textbf{La risposta di Cutolo:} In algebra e teoria degli insiemi, i numeri naturali rappresentano le \textbf{cardinalità degli insiemi finiti}. L'insieme vuoto $\vuoto$ è un insieme finito con zero elementi, dunque $|\vuoto| = 0 \in \N$.
Inoltre, negli assiomi di Peano formulati da Dedekind, l'elemento base può essere $0$, consentendo a $(\N, +, 0)$ di essere un monoide commutativo con elemento neutro $0$.

\subsection{2. Perché si scrive $\antivecf$ e MAI $f^{-1}$ se $f$ non è biettiva?}
\textbf{La trappola:} Usare la notazione $f^{-1}(B)$ per indicare la controimmagine di un sottoinsieme.
\textbf{La risposta di Cutolo:} La scrittura $f^{-1}$ denota l'\textbf{applicazione inversa} $f^{-1}: B \to A$, la quale esiste \textbf{se e solo se $f$ è biettiva}.
L'operazione di controimmagine $\antivecf: \P(B) \to \P(A)$ definita da $\antivecf(Y) = \{x \in A \mid f(x) \in Y\}$ è invece un'applicazione perfettamente definita per \textbf{qualunque} applicazione $f$, anche quando $f$ non è né iniettiva né suriettiva. Per evitare l'errore concettuale di presupporre l'invertibilità di $f$, si usa la freccia indietro $\antivecf$.

\subsection{3. Dimostrare che $\inf$ è associativo senza usare gli assiomi algebrici}
\textbf{La trappola:} Recitare l'uguaglianza $(a \wedge b) \wedge c = a \wedge (b \wedge c)$ dicendo che è un assioma del reticolo.
\textbf{La risposta di Cutolo:} Si deve lavorare unicamente nella definizione d'ordine:
\begin{enumerate}
    \item Si definisce $\inf(S) = \max(\minor(S))$.
    \item Si mostra che sia $d_1 = a \wedge (b \wedge c)$ che $d_2 = (a \wedge b) \wedge c$ sono entrambi il massimo dell'insieme dei minoranti comuni della terna $\{a, b, c\}$.
    \item Per l'unicità del massimo, $d_1 = d_2$. (Si veda la dimostrazione formale nel Teorema~\ref{teo:teo_equiv_reticolo}).
\end{enumerate}

\subsection{4. Cancellabilità vs Invertibilità: Perché in un monoide finito coincidono?}
\textbf{La trappola:} Dare per scontato che cancellabile $\implies$ invertibile.
\textbf{La risposta di Cutolo:}
\begin{itemize}
    \item In un monoide qualsiasi, invertibile $\implies$ cancellabile (si moltiplica per l'inverso).
    \item Il viceversa è falso nei monoidi infiniti: in $(\Z, \cdot)$, $2$ è cancellabile ($2x = 2y \implies x = y$), ma non ammette inverso in $\Z$.
    \item Se il monoide $M$ è \textbf{finito}, la cancellabilità a sinistra di $a$ equivale all'\textbf{iniettività} della traslazione $\lambda_a: M \to M$. Poiché $M$ è finito, ogni applicazione iniettiva da $M$ in $se$ stesso è anche \textbf{suriettiva}. Dunque esiste $b \in M$ tale che $\lambda_a(b) = ab = 1$, provando che $a$ è invertibile.
\end{itemize}

\subsection{5. Cancellabilità e Divisori dello Zero negli Anelli}
\textbf{La domanda:} Qual è l'esatta relazione tra elementi cancellabili e divisori dello zero in un anello commutativo unitario $R$?
\textbf{La dimostrazione:}
Un elemento $a \in R$ è cancellabile rispetto alla moltiplicazione se e solo se $a$ \textbf{non è un divisore dello zero}:
\begin{itemize}
    \item $(\implica)$ Sia $a$ cancellabile e supponiamo $ab = 0$. Poiché $a \cdot 0 = 0$, abbiamo $ab = a \cdot 0$. Cancellando $a$ otteniamo $b = 0$. Quindi $a$ non può essere un divisore dello zero.
    \item $(\impliedby)$ Sia $a$ non divisore dello zero e supponiamo $ax = ay$. Allora $ax - ay = 0 \implies a(x - y) = 0$. Poiché $a$ non è divisore dello zero e $a \neq 0$, deve necessariamente essere $x - y = 0$, da cui $x = y$. Dunque $a$ è cancellabile.
\end{itemize}

\subsection{6. La Dimostrazione del Teorema di Ruffini dal Teorema del Resto}
\textbf{La trappola:} Spiegare la regola tabellare di Ruffini per calcolare i coefficienti.
\textbf{La risposta di Cutolo:} Il Teorema di Ruffini afferma che $c$ è radice di $f(x) \iff (x - c) \mid f(x)$.
La dimostrazione discende dalla divisione euclidea: $f(x) = (x - c)q(x) + r(x)$ con $\degpoly(r) < 1$, quindi $r \in K$ è una costante.
Valutando l'omomorfismo di valutazione in $x = c$: $f(c) = (c - c)q(c) + r = 0 + r = r$.
Dunque il resto della divisione è esattamente $f(c)$. Ne segue che $(x - c) \mid f(x) \iff \text{resto nullo } r = 0 \iff f(c) = 0$.

\subsection{7. Perché il Criterio delle Radici fallisce per Gradi Superiori a 3?}
\textbf{La trappola:} Dire che se un polinomio non ha radici è sempre irriducibile.
\textbf{La risposta di Cutolo:} Un polinomio di secondo o terzo grado, se fosse riducibile, dovrebbe avere almeno un fattore di primo grado (poiché la somma dei gradi è 2 o 3), e tale fattore lineare fornirebbe una radice.
Per un polinomio di quarto grado, invece, la riducibilità può avvenire come prodotto di due polinomi irriducibili di secondo grado:
\[
    f(x) = g(x) \cdot h(x) \qquad \text{con } \degpoly(g) = \degpoly(h) = 2
\]
In tal caso $f(x)$ \textbf{non ammette alcuna radice nel campo}, ma è comunque \textbf{riducibile}!

\subsection{8. Sottoreticolo vs Sottoinsieme che è un Reticolo}
\textbf{La trappola:} Credere che ogni sottoinsieme di un reticolo che sia un reticolo con l'ordine indotto sia un sottoreticolo.
\textbf{La risposta di Cutolo:} Un sottoinsieme $S \sseq L$ è un \textbf{sottoreticolo} se è \textbf{chiuso rispetto alle operazioni} del reticolo ambiente:
\[
    \forall a, b \in S, \quad a \wedge_L b \in S \;\wedge\; a \vee_L b \in S
\]
Può accadere che $S$, considerato come poset a sé stante con l'ordine indotto, sia un reticolo, ma che l'estremo inferiore di due elementi calcolato in $S$ \textbf{non coincida} con l'estremo inferiore calcolato nel reticolo ambiente $L$! In tal caso $S$ non è un sottoreticolo di $L$.

\subsection{9. Perché l'Anello Quoziente $K[x]/(f(x))$ è un Campo solo se $f$ è Irriducibile?}
\textbf{La spiegazione:}
\begin{itemize}
    \item Se $f(x) = g(x)h(x)$ è riducibile, allora $[g] \cdot [h] = [f] = [0]$ con $[g], [h] \neq [0]$, dunque il quoziente ha divisori dello zero e non può essere un campo.
    \item Se $f(x)$ è irriducibile, per ogni $[g] \neq [0]$ si ha $\mcd(f, g) = 1$. Per il Teorema di Bézout per polinomi, $f(x)u(x) + g(x)v(x) = 1$. Passando alle classi resto: $[g(x)][v(x)] = [1]$, dunque ogni elemento non nullo è invertibile, rendendo il quoziente un campo.
\end{itemize}

\subsection{10. Disegnare i Reticoli di Birkhoff e Indicare la Violazione}
\textbf{La richiesta:} Disegnare al volo i diagrammi di Hasse del pentagono $N_5$ e del diamante $M_3$ ed esibire la terna che non rispetta la distributività. (Si veda la Figura~\ref{fig:n5_m3} nel Capitolo~\ref{chap:ordini_reticoli}).

\section{I Dieci Comandamenti del Rigore Matematico per l'Orale}

\begin{enumerate}
    \item \textbf{Mai omettere i quantificatori:} Scrivere sempre $\forall$ o $\exists$ specificando l'insieme di appartenenza.
    \item \textbf{Dichiarare le variabili prima di usarle:} Non iniziare una frase con ``$f(x) = y$'' senza aver prima dichiarato ``Siano $A, B$ insiemi e sia $f: A \to B$ un'applicazione''.
    \item \textbf{Distinguere implicazione logica da equivalenza:} Non usare $\iff$ al posto di $\implies$.
    \item \textbf{Non confondere appartenenza ($\in$) ed inclusione ($\sseq$):} $a \in A$ connette un elemento a un insieme; $\{a\} \sseq A$ connette due insiemi.
    \item \textbf{Non dire ``si annulla'' in modo ambiguo:} Specificare se ci si riferisce a un elemento uguale a zero ($a = 0$) o a un elemento che è divisore dello zero.
    \item \textbf{Enunciare sempre le ipotesi prima della tesi:} Nei teoremi le ipotesi non sono dettagli opzionali; se si omette l'ipotesi di campo per la divisione euclidea dei polinomi, il teorema è falso.
    \item \textbf{Rispettare le convenzioni del docente:} Usare $\antivecf$ per le controimmagini e includere $0$ in $\N$.
    \item \textbf{Nei reticoli, motivare $\inf$ e $\sup$ con la definizione:} Non limitarsi a dire ``si vede dal disegno''.
    \item \textbf{Conoscere a memoria gli enunciati esatti:} All'orale Cutolo chiede spesso: ``Mi enunci il Teorema di Lagrange'', ``Mi enunci il Teorema di Stone'', ``Mi enunci il Criterio di Birkhoff''.
    \item \textbf{Se si commette un errore, fermarsi e correggersi con calma:} Cutolo apprezza grandemente lo studente che, resosi conto di un'imprecisione, torna indietro, analizza l'assioma violato e riformula l'argomentazione con rigore formale.
\end{enumerate}

\chapter{Eserciziario Risolto e Commentato Tipo-Esame}
\label{chap:eserciziario}

In questo capitolo presentiamo una selezione di esercizi d'esame accuratamente risolti e commentati, rappresentativi di tutte le tipologie di problemi assegnati nelle prove scritte del Prof. Cutolo. Per ciascun esercizio viene evidenziata la metodologia risolutiva generale e i passaggi formali attesi in sede di correzione.

% ------------------------------------------------------------------------------
\section{Studio di Relazioni di Equivalenza e Quozienti}

\begin{esercizio}{Relazione d'Equivalenza su $\Z \times \Z$}{es_equiv}
Sull'insieme $A = \Z \times \Z$ si consideri la relazione binaria $\mathcal{R}$ definita ponendo, per ogni $(a, b), (c, d) \in A$:
\[
    (a, b) \mathcal{R} (c, d) \quad \Shiff \quad a + 2d = c + 2b
\]
\begin{enumerate}
    \item Dimostrare che $\mathcal{R}$ è una relazione di equivalenza su $A$.
    \item Determinare esplicitamente la classe di equivalenza $[(0, 0)]_\mathcal{R}$ e la classe $[(1, 2)]_\mathcal{R}$.
    \item Descrivere geometricamente o algebricamente l'insieme quoziente $A/\mathcal{R}$.
\end{enumerate}
\end{esercizio}

\begin{proof}[Risoluzione Commentata]
\textbf{Punto 1: Verifica delle proprietà di equivalenza.}
Riscriviamo la condizione spostando le variabili con lo stesso indice:
\[
    a + 2d = c + 2b \quad \Shiff \quad a - 2b = c - 2d
\]
Definiamo l'applicazione ausiliaria $f: \Z \times \Z \to \Z$ ponendo $f(a, b) = a - 2b$.
Allora la relazione $\mathcal{R}$ coincide esattamente con il \textbf{nucleo di equivalenza} di $f$:
\[
    (a, b) \mathcal{R} (c, d) \quad \Shiff \quad f(a, b) = f(c, d)
\]
Poiché il nucleo di equivalenza di qualunque applicazione è una relazione di equivalenza (Definizione~\ref{def:nucleo_eq}), $\mathcal{R}$ è automaticamente riflessiva, simmetrica e transitiva!
\begin{itemize}
    \item \emph{Riflessività:} $\forall (a, b)$, $f(a, b) = f(a, b) \implies a - 2b = a - 2b \implies (a, b) \mathcal{R} (a, b)$.
    \item \emph{Simmetria:} se $f(a, b) = f(c, d)$, allora $f(c, d) = f(a, b)$, quindi $(c, d) \mathcal{R} (a, b)$.
    \item \emph{Transitività:} se $f(a, b) = f(c, d)$ e $f(c, d) = f(e, g)$, per transitività dell'uguaglianza in $\Z$ si ha $f(a, b) = f(e, g)$, da cui $(a, b) \mathcal{R} (e, g)$.
\end{itemize}

\textbf{Punto 2: Determinazione delle classi di equivalenza.}
Per definizione, la classe di $(x_0, y_0)$ è la fibra $\antivecf(\{f(x_0, y_0)\})$:
\begin{itemize}
    \item Per $(0, 0)$: $f(0, 0) = 0 - 0 = 0$. Dunque:
    \[
        [(0, 0)]_\mathcal{R} = \{(a, b) \in \Z \times \Z \mid a - 2b = 0\} = \{(2b, b) \mid b \in \Z\}
    \]
    È l'insieme di tutti i punti a coordinate intere giacenti sulla retta $a = 2b$.
    \item Per $(1, 2)$: $f(1, 2) = 1 - 2(2) = -3$. Dunque:
    \[
        [(1, 2)]_\mathcal{R} = \{(a, b) \in \Z \times \Z \mid a - 2b = -3\} = \{(2b - 3, b) \mid b \in \Z\}
    \]
\end{itemize}

\textbf{Punto 3: Struttura dell'insieme quoziente $A/\mathcal{R}$.}
Per il Teorema di Omomorfismo per insiemi (Teorema~\ref{teo:teo_omom_insiemi}), l'applicazione indotta:
\[
    \tilde{f}: (\Z \times \Z)/\mathcal{R} \longrightarrow \im(f), \qquad [(a, b)]_\mathcal{R} \longmapsto a - 2b
\]
è una \textbf{biiezione}.
Poiché per ogni $k \in \Z$ si ha $f(k, 0) = k - 0 = k$, l'applicazione $f$ è suriettiva su $\Z$, quindi $\im(f) = \Z$.
Ne deduciamo che l'insieme quoziente $(\Z \times \Z)/\mathcal{R}$ è in biiezione con $\Z$ (ha la cardinalità del numerabile $\aleph_0$), e ogni classe corrisponde a una retta discreta di pendenza $1/2$ nel reticolo $\Z^2$.
\end{proof}

% ------------------------------------------------------------------------------
\section{Studio di un Poset e Diagramma di Hasse}

\begin{esercizio}{Poset dei Divisori e Verifica di Reticolo}{es_poset}
Sia $S = \{2, 3, 4, 6, 8, 12, 24\}$ ordinato mediante la relazione di divisibilità $\mid$ ($a \le b \iff a \mid b$).
\begin{enumerate}
    \item Verificare che $(S, \mid)$ è un poset e tracciarne il diagramma di Hasse.
    \item Determinare gli elementi minimali, massimali, l'eventuale minimo e massimo di $S$.
    \item Stabilire se $(S, \mid)$ è un reticolo.
    \item Calcolare $\inf(4, 6)$ e $\sup(4, 6)$ in $S$.
\end{enumerate}
\end{esercizio}

\begin{proof}[Risoluzione Commentata]
\textbf{Punto 1: Diagramma di Hasse.}
La relazione di divisibilità su interi positivi è un ordine parziale (riflessiva, antisimmetrica, transitiva).
Troviamo le relazioni di copertura $x \prec y$ ($x \mid y$ senza divisori intermedi in $S$):
\begin{itemize}
    \item Da $2$: $2 \prec 4$ (non ci sono intermedi), $2 \prec 6$. (Non copre $8$ perché $2 \mid 4 \mid 8$).
    \item Da $3$: $3 \prec 6$.
    \item Da $4$: $4 \prec 8$, $4 \prec 12$.
    \item Da $6$: $6 \prec 12$.
    \item Da $8$: $8 \prec 24$.
    \item Da $12$: $12 \prec 24$.
\end{itemize}

Il diagramma di Hasse è illustrato di seguito:

\begin{center}
\begin{tikzpicture}[scale=1.2, every node/.style={circle, draw, fill=white, inner sep=2.5pt}]
    \node (2) at (-1, 0) {$2$};
    \node (3) at (1, 0) {$3$};
    \node (4) at (-1.5, 1.3) {$4$};
    \node (6) at (0.2, 1.3) {$6$};
    \node (8) at (-1.5, 2.6) {$8$};
    \node (12) at (0.2, 2.6) {$12$};
    \node (24) at (-0.6, 3.8) {$24$};
    
    \draw (2) -- (4);
    \draw (2) -- (6);
    \draw (3) -- (6);
    \draw (4) -- (8);
    \draw (4) -- (12);
    \draw (6) -- (12);
    \draw (8) -- (24);
    \draw (12) -- (24);
\end{tikzpicture}
\end{center}

\textbf{Punto 2: Elementi Speciali.}
\begin{itemize}
    \item \textbf{Minimali:} gli elementi che non hanno nulla al di sotto: $\{2, 3\}$.
    \item \textbf{Minimo:} non esiste alcun elemento che divida tutti gli elementi di $S$ contemporaneamente ($2 \nmid 3$ e $3 \nmid 2$). Dunque \textbf{non esiste minimo}.
    \item \textbf{Massimali:} l'unico elemento che non ha nulla al di sopra è $\{24\}$.
    \item \textbf{Massimo:} poiché $24$ è multiplo di tutti gli elementi di $S$ ($2, 3, 4, 6, 8, 12$ dividono tutti $24$), \textbf{$24$ è il massimo assoluto}: $\max(S) = 24$.
\end{itemize}

\textbf{Punti 3 e 4: Verifica di Reticolo.}
Per essere un reticolo, ogni coppia di elementi deve ammettere $\inf$ e $\sup$ \textbf{in $S$}.
Consideriamo la coppia $\{2, 3\}$:
I minoranti di $\{2, 3\}$ in $S$ sono: $\minor(\{2, 3\}) = \{x \in S \mid x \mid 2 \wedge x \mid 3\} = \vuoto$ (poiché l'unico divisore comune sarebbe $1$, ma $1 \notin S$).
Non essendovi minoranti in $S$, l'insieme $\minor(\{2, 3\})$ non può avere massimo, dunque \textbf{non esiste $\inf(2, 3)$ in $S$}.
Ne segue immediatamente che \textbf{$(S, \mid)$ non è un reticolo}!

Calcoliamo ora $\inf(4, 6)$ e $\sup(4, 6)$:
\begin{itemize}
    \item Minoranti di $\{4, 6\}$ in $S$: $\minor(\{4, 6\}) = \{x \in S \mid x \mid 4 \wedge x \mid 6\} = \{2\}$.
    Il massimo di tale insieme è $2$, quindi:
    \[
        \inf(4, 6) = 2
    \]
    \item Maggioranti di $\{4, 6\}$ in $S$: $\maggior(\{4, 6\}) = \{x \in S \mid 4 \mid x \wedge 6 \mid x\} = \{12, 24\}$.
    Il minimo di questo insieme è $12$, dunque:
    \[
        \sup(4, 6) = 12
    \]
\end{itemize}
\end{proof}

% ------------------------------------------------------------------------------
\section{Equazione Congruenziale Lineare con Algoritmo di Bézout}

\begin{esercizio}{Risoluzione Completa di Congruenza Lineare}{es_congr}
Si consideri l'equazione congruenziale lineare:
\[
    72x \equiv 42 \pmod{102}
\]
\begin{enumerate}
    \item Stabilire se l'equazione ammette soluzioni e quante soluzioni distinte modulo $102$ esistono.
    \item Determinare tutte le soluzioni intere in $\Z$ e tutte le soluzioni modulo $102$.
\end{enumerate}
\end{esercizio}

\begin{proof}[Risoluzione Commentata]
\textbf{Passo 1: Calcolo di $d = \mcd(72, 102)$ con l'algoritmo euclideo.}
Eseguiamo le divisioni successive:
\begin{align*}
    102 &= 72 \cdot 1 + 30 \\
    72 &= 30 \cdot 2 + 12 \\
    30 &= 12 \cdot 2 + 6 \\
    12 &= 6 \cdot 2 + 0 \quad (\text{resto nullo})
\end{align*}
L'ultimo resto non nullo è $d = 6$. Dunque $\mcd(72, 102) = 6$.
Poiché $d = 6$ divide il termine noto $42$ ($42 = 6 \cdot 7$), per il Criterio di Solvibilità (Teorema~\ref{teo:solv_congr}) l'equazione \textbf{è risolubile} e ammette esattamente \textbf{$6$ soluzioni distinte modulo $102$}.

\textbf{Passo 2: Semplificazione dell'equazione.}
Dividiamo tutti i termini per $d = 6$:
\[
    \frac{72}{6} x \equiv \frac{42}{6} \pmod{\frac{102}{6}} \quad \Implica \quad 12x \equiv 7 \pmod{17}
\]
I coefficienti $12$ e $17$ sono coprimi: $\mcd(12, 17) = 1$.

\textbf{Passo 3: Calcolo dell'inverso di $12$ modulo $17$ con l'identità di Bézout.}
Eseguiamo le divisioni tra $17$ e $12$:
\begin{align*}
    17 &= 12 \cdot 1 + 5 \\
    12 &= 5 \cdot 2 + 2 \\
    5 &= 2 \cdot 2 + 1
\end{align*}
Risaliamo a ritroso per esprimere $1$ come combinazione lineare di $17$ e $12$:
\begin{align*}
    1 &= 5 - 2 \cdot 2 \\
    &= 5 - 2 \cdot (12 - 5 \cdot 2) = 5 \cdot 5 - 2 \cdot 12 \\
    &= 5 \cdot (17 - 12 \cdot 1) - 2 \cdot 12 = 5 \cdot 17 - 7 \cdot 12
\end{align*}
Abbiamo ottenuto l'identità di Bézout:
\[
    17 \cdot 5 + 12 \cdot (-7) = 1
\]
Passando modulo $17$:
\[
    12 \cdot (-7) \equiv 1 \pmod{17}
\]
Poiché $-7 \equiv 10 \pmod{17}$, l'inverso di $12$ modulo $17$ è $[10]_{17}$.

\textbf{Passo 4: Soluzione particolare.}
Moltiplichiamo entrambi i membri di $12x \equiv 7 \pmod{17}$ per l'inverso $10$:
\[
    x_0 \equiv 10 \cdot 7 = 70 \equiv 2 \pmod{17} \quad (\text{poiché } 70 = 17 \cdot 4 + 2)
\]
Dunque una soluzione particolare è $x_0 = 2$.

\textbf{Passo 5: Tutte le $6$ soluzioni modulo $102$.}
Le $d = 6$ soluzioni distinte modulo $102$ si ottengono aggiungendo multipli di $m' = 17$:
\[
    x_k = x_0 + k \cdot 17 \qquad \text{per } k \in \{0, 1, 2, 3, 4, 5\}
\]
Calcoliamo esplicitamente:
\begin{align*}
    x_0 &= 2 \\
    x_1 &= 2 + 17 = 19 \\
    x_2 &= 19 + 17 = 36 \\
    x_3 &= 36 + 17 = 53 \\
    x_4 &= 53 + 17 = 70 \\
    x_5 &= 70 + 17 = 87
\end{align*}
Le classi di equivalenza delle soluzioni in $\Z_{102}$ sono $\{[2], [19], [36], [53], [70], [87]\}$.
L'insieme di tutte le soluzioni intere in $\Z$ è:
\[
    S = \{2 + 17k \mid k \in \Z\}
\]
\end{proof}

% ------------------------------------------------------------------------------
\section{Calcolo di Grandi Potenze Modulari}

\begin{esercizio}{Calcolo Potenze con Eulero-Fermat}{es_potenza}
Determinare il resto della divisione di $13^{2026}$ per $35$.
\end{esercizio}

\begin{proof}[Risoluzione Commentata]
Vogliamo calcolare $13^{2026} \pmod{35}$.
\begin{enumerate}
    \item \textbf{Verifica della coprimarità:} La fattorizzazione in primi di $35$ è $35 = 5 \cdot 7$.
    Poiché $13$ non è divisibile né per $5$ né per $7$, si ha $\mcd(13, 35) = 1$.
    Possiamo quindi applicare il \textbf{Teorema di Eulero-Fermat}.
    \item \textbf{Calcolo di $\varphi(35)$:} Per la proprietà moltiplicativa:
    \[
        \varphi(35) = \varphi(5) \cdot \varphi(7) = (5 - 1) \cdot (7 - 1) = 4 \cdot 6 = 24
    \]
    Per il teorema: $13^{24} \equiv 1 \pmod{35}$.
    \item \textbf{Divisione dell'esponente per $\varphi(35) = 24$:}
    \[
        2026 = 24 \cdot 84 + 10 \qquad (\text{resto } 10)
    \]
    \item \textbf{Riduzione dell'esponente:}
    \[
        13^{2026} = (13^{24})^{84} \cdot 13^{10} \equiv 1^{84} \cdot 13^{10} \equiv 13^{10} \pmod{35}
    \]
    \item \textbf{Calcolo rapido di $13^{10} \pmod{35}$:}
    \begin{align*}
        13^2 &= 169 = 35 \cdot 4 + 29 \equiv -6 \pmod{35} \\
        13^4 &= (13^2)^2 \equiv (-6)^2 = 36 \equiv 1 \pmod{35} \\
        13^8 &= (13^4)^2 \equiv 1^2 = 1 \pmod{35} \\
        13^{10} &= 13^8 \cdot 13^2 \equiv 1 \cdot (-6) = -6 \equiv 29 \pmod{35}
    \end{align*}
\end{enumerate}
Il resto cercato è quindi $29$.
\end{proof}

% ------------------------------------------------------------------------------
\section{Fattorizzazione di Polinomi in $\Z_p[x]$ e Campi Quoziente}

\begin{esercizio}{Irriducibilità e Quoziente in $\Z_3[x]$}{es_polinomio}
In $\Z_3[x]$, si consideri il polinomio $f(x) = x^3 + 2x + 1$.
\begin{enumerate}
    \item Determinare se $f(x)$ è irriducibile in $\Z_3[x]$.
    \item Stabilire se l'anello quoziente $K = \Z_3[x]/(f(x))$ è un campo e calcolarne la cardinalità $|K|$.
    \item Nel quoziente $K$, determinare l'inverso moltiplicativo dell'elemento $[x]$.
\end{enumerate}
\end{esercizio}

\begin{proof}[Risoluzione Commentata]
\textbf{Punto 1: Studio dell'irriducibilità.}
Il polinomio $f(x)$ ha grado $3$. Per la Proposizione~\ref{prop:criterio_grado_2_3}, un polinomio di terzo grado a coefficienti in un campo è irriducibile se e solo se non possiede radici nel campo.
Valutiamo $f(x)$ in tutti gli elementi di $\Z_3 = \{0, 1, 2\}$:
\begin{align*}
    f(0) &= 0^3 + 2(0) + 1 = 1 \not\equiv 0 \pmod 3 \\
    f(1) &= 1^3 + 2(1) + 1 = 1 + 2 + 1 = 4 \equiv 1 \not\equiv 0 \pmod 3 \\
    f(2) &= 2^3 + 2(2) + 1 = 8 + 4 + 1 = 13 \equiv 1 \not\equiv 0 \pmod 3
\end{align*}
Poiché $f(c) \neq 0$ per ogni $c \in \Z_3$, $f(x)$ \textbf{non ha radici in $\Z_3$}.
Essendo di terzo grado, ne deduciamo che \textbf{$f(x)$ è irriducibile in $\Z_3[x]$}.

\textbf{Punto 2: Natura e cardinalità del quoziente.}
Per il Teorema~\ref{teo:campo_quoziente}, poiché $f(x)$ è irriducibile su un campo, l'anello quoziente:
\[
    K = \Z_3[x]/(f(x))
\]
è un \textbf{campo} (noto come campo finito o estensione di Galois $\F_{3^3} = \F_{27}$).
I suoi elementi sono rappresentati univocamente dai polinomi di grado strettamente minore di $\degpoly(f) = 3$:
\[
    K = \{a_0 + a_1 x + a_2 x^2 + (f(x)) \mid a_0, a_1, a_2 \in \Z_3\}
\]
Poiché ci sono 3 scelte indipendenti per ciascuno dei 3 coefficienti ($a_0, a_1, a_2$), la cardinalità è:
\[
    |K| = 3^3 = 27 \text{ elementi}
\]

\textbf{Punto 3: Calcolo dell'inverso di $[x]$ in $K$.}
Nel quoziente vale la relazione $[f(x)] = [0]$, cioè:
\[
    [x^3 + 2x + 1] = [0] \quad \Implica \quad [x^3 + 2x] = [-1] = [2]
\]
Raccogliendo $[x]$ a fattor comune a sinistra:
\[
    [x] \cdot [x^2 + 2] = [2]
\]
Moltiplichiamo entrambi i membri per l'inverso di $2$ in $\Z_3$ (che è $2$, poiché $2 \cdot 2 = 4 \equiv 1 \pmod 3$):
\[
    [x] \cdot [2(x^2 + 2)] = [2 \cdot 2] = [1]
\]
Poiché $2(x^2 + 2) = 2x^2 + 4 \equiv 2x^2 + 1 \pmod 3$, otteniamo:
\[
    [x]^{-1} = [2x^2 + 1]
\]
Verifica immediata: $x(2x^2 + 1) = 2x^3 + x = 2(-2x - 1) + x = -4x - 2 + x = -3x - 2 \equiv -2 \equiv 1 \pmod 3$.
\end{proof}


\end{document}
